

\begin{register}{H}{DMABUSMODE\_REG}{0x1000}\label{regdesc:DMABUSMODEREG}
\regfield{(DMAREBUILDINCRBURST)}{1}{31}{0}%
\regfield{(reserved)}{3}{28}{00000}%
\regfield{DMATRANSMITPRIORITY}{1}{27}{0}%
\regfield{DMAMIXEDBURST}{1}{26}{0}%
\regfield{DMAADDRALIBEA}{1}{25}{0}%
\regfield{PBLX8\_MODE}{1}{24}{0}%
\regfield{USE\_SEP\_PBL}{1}{23}{0}%
\regfield{RX\_DMA\_PBL}{6}{17}{{0x01}}%
\regfield{FIXED\_BURST}{1}{16}{0}%
\regfield{PRI\_RATIO}{2}{14}{{0x0}}%
\regfield{PROG\_BURST\_LEN}{6}{8}{{0x01}}%
\regfield{ALT\_DESC\_SIZE}{1}{7}{0}%
\regfield{DESC\_SKIP\_LEN}{5}{2}{{0x00}}%
\regfield{DMA\_ARB\_SCH}{1}{1}{0}%
\regfield{SW\_RST}{1}{0}{1}%
\reglabel{Reset}\regnewline%
\begin{regdesc}\begin{reglist}
\item [DMAREBUILDINCRBURST] Configures whether to rebuild INCRx burst.\\
0: Not rebuild\\
1: If the AHB master gets an EBT (retry, split, or losing bus grant), the AHB master interface rebuilds the pending beats of any burst transfer initiated with INCRx using INCRx and SINGLE.\\(R/W)
\item [DMATRANSMITPRIORITY] Configures whether the transmit DMA has higher priority than the receive DMA during arbitration.\\
0: Lower priority\\
1: Higher priority\\ (R/W)
\item [DMAMIXEDBURST] Configures whether to use mixed burst.\\
0: Not use\\
1: If the FIXED\_BURST bit is 1, the AHB Master interface starts all bursts of length more than 16 with INCR (undefined burst) whereas it reverts to fixed burst transfers (INCRx and SINGLE) for burst length of 16 and less.\\(R/W)
\item [DMAADDRALIBEA] Configures whether bursts are aligned to address.\\
0: Not aligned to address\\
1: If the FIXED\_BURST bit is 1, the AHB interface generates all bursts aligned to the start address LS bits. If the FIXED\_BURST bit is 0, the first burst (accessing the data buffer's start address) is not aligned, but subsequent bursts are aligned to the address.\\(R/W)

\item [Continued on the next page...]
\end{reglist}\end{regdesc}
\end{register}

\addtocounter{Regfloat}{-1}
\begin{register}{H}{DMABUSMODE\_REG}{0x1000}
\begin{regdesc}\begin{reglist}
\item [Continued from the previous page...]

\item [PBLX8\_MODE] Configures whether to multiply the PBL (programmable burst length) value by 8.\\
0: Disable\\
1: The programmed PBL value (PROG\_BURST\_LEN and RX\_DMA\_PBL) is multiplied by 8. Therefore, the DMA transfers the data in 8, 16, 32, 64, 128, and 256 beats depending on the PBL value.\\ (R/W)
\item [USE\_SEP\_PBL] Configures whether to use separate PBL.\\
0: PROG\_BURST\_LEN is applicable only for TX DMA, and RX DMA uses RX\_DMA\_PBL.\\
1: PROG\_BURST\_LEN is applicable for both TX and RX DMA.\\(R/W)
\item [RX\_DMA\_PBL] Configure the maximum number of beats to be transferred in one RX DMA transaction, namely the maximum value that is used in a single block Read or Write.\\
Value range: 1, 2, 4, 8, 16 and 32. Any other value results in undefined behavior.\\
Valid only when USE\_SEP\_PBL is set.\\
(R/W)
\item [FIXED\_BURST] Configures whether the AHB Master interface performs fixed burst transfers.\\
0: The AHB interface uses SINGLE and INCR burst transfer operations.\\
1: The AHB interface uses only SINGLE, INCR4, INCR8, or INCR16 during start of the normal burst transfers.\\ (R/W)
\item [PRI\_RATIO] Configures the priority ratio in the weighted round-robin arbitration between the RX DMA and TX DMA.\\
0: 1:1\\
1: 2:0\\
2: 3:1\\
3: 4:1\\
Valid only when DMA\_ARB\_SCH is 0. (R/W)
\item [PROG\_BURST\_LEN]
Configures the maximum number of beats to be transferred in one DMA transaction.\\
If the number of beats to be transferred is more than 32, then perform the following steps: \\
1. Set the PBLx8 mode with PBLX8\_MODE\\
2. Set the PBL with PROG\_BURST\_LEN\\(R/W)
\item [ALT\_DESC\_SIZE] Configures whether to increase the descriptor size to 32 bytes.\\
0: Keep 16 bytes\\
1: Increase to 32 bytes\\(R/W)
\item [DESC\_SKIP\_LEN] Configures the number of Word, Dword, or Lword (depending on the 32-bit, 64-bit, or 128-bit bus) to skip between two unchained descriptors.\\The address skipping starts from the end of the current descriptor to the start of next descriptor. When this field is 0, the descriptor table is taken as contiguous by the DMA in Ring mode. (R/W)

\item [Continued on the next page...]
\end{reglist}\end{regdesc}
\end{register}

\addtocounter{Regfloat}{-1}
\begin{register}{H}{DMABUSMODE\_REG}{0x1000}
\begin{regdesc}\begin{reglist}
\item [Continued from the previous page...]

\item [DMA\_ARB\_SCH] Configures the arbitration scheme between the transmit and receive paths.\\
0: Weighted round-robin with RX:TX or TX:RX. In this case, the priority between the paths is according to the priority specified in PRI\_RATIO.\\
1: Fixed priority (The transmit path has priority over receive path).\\(R/W)
\item [SW\_RST] Configures whether the MAC DMA Controller resets the logic and all internal registers of the MAC.\\
0: Release from reset\\
1: Reset\\
This bit is cleared automatically after the reset operation has completed in all of the ETH\_MAC clock domains.\\ (R/WS/SC)
\end{reglist}\end{regdesc}
\end{register}


\begin{register}{H}{DMATXPOLLDEMAND\_REG}{0x1004}\label{regdesc:DMATXPOLLDEMANDREG}
\regfield{TRANS\_POLL\_DEMAND}{32}{0}{{0x000000000}}%
\reglabel{Reset}\regnewline%
\begin{regdesc}\begin{reglist}
\item [TRANS\_POLL\_DEMAND] Configures whether to enable the TX DMA to check if the DMA owns the current descriptor. \\
Any value: Enable\\
When this field is written, the DMA reads the current descriptor pointed to by \hyperref[regdesc:DMATXCURRDESCREG]{DMATXCURRDESC\_REG}. If that descriptor is not available (owned by the Host), the transmission returns to the Suspend state and the Bit 2 (TU) of Register 5 (Status Register) is asserted. If the descriptor is available, the transmission resumes. (RO/WT)
\end{reglist}\end{regdesc}
\end{register}


\begin{register}{H}{DMARXPOLLDEMAND\_REG}{0x1008}\label{regdesc:DMARXPOLLDEMANDREG}
\regfield{RECV\_POLL\_DEMAND}{32}{0}{{0x000000000}}%
\reglabel{Reset}\regnewline%
\begin{regdesc}\begin{reglist}
\item [RECV\_POLL\_DEMAND] Configures whether to enable the RX DMA to check if the DMA owns the current descriptor. \\
Any value: Enable\\
When this field is written, the DMA reads the current descriptor pointed to by \hyperref[regdesc:DMARXCURRDESCREG]{DMARXCURRDESC\_REG}. If that descriptor is not available (owned by the Host), the reception returns to the Suspended state and the Bit 7 (RU) of Register 5 (Status Register) is not asserted. If the descriptor is available, the RX DMA returns to the active state.(RO/WT)
\end{reglist}\end{regdesc}
\end{register}


\begin{register}{H}{DMARXBASEADDR\_REG}{0x100C}\label{regdesc:DMARXBASEADDRREG}
\regfield{START\_RECV\_LIST}{32}{0}{{0x000000000}}%
\reglabel{Reset}\regnewline%
\begin{regdesc}\begin{reglist}
\item [START\_RECV\_LIST] Configures the base address of the first descriptor in the receive descriptor linked list.\\ The LSB bits (1:0) are ignored (32-bit wide bus) and internally taken as all-zero by the DMA. Therefore, these LSB bits are read-only (RO). (R/W)
\end{reglist}\end{regdesc}
\end{register}


\begin{register}{H}{DMATXBASEADDR\_REG}{0x1010}\label{regdesc:DMATXBASEADDRREG}
\regfield{START\_TRANS\_LIST}{32}{0}{{0x000000000}}%
\reglabel{Reset}\regnewline%
\begin{regdesc}\begin{reglist}
\item [START\_TRANS\_LIST] Configures the base address of the first descriptor in the Transmit Descriptor list. The LSB bits (1:0) are ignored (32-bit wide bus) and are internally taken as all-zero by the DMA. Therefore, these LSB bits are read-only (RO). (R/W)
\end{reglist}\end{regdesc}
\end{register}


\begin{register}{H}{DMASTATUS\_REG}{0x1014}\label{regdesc:DMASTATUSREG}
\regfield{(reserved)}{1}{31}{0}%
\regfield{EMAC\_LPI\_INT}{1}{30}{0}%
\regfield{TS\_TRI\_INT}{1}{29}{0}%
\regfield{EMAC\_PMT\_INT}{1}{28}{0}%
\regfield{(reserved)}{2}{26}{00}%
\regfield{ERROR\_BITS}{3}{23}{{0x0}}%
\regfield{TRANS\_PROC\_STATE}{3}{20}{{0x0}}%
\regfield{RECV\_PROC\_STATE}{3}{17}{{0x0}}%
\regfield{NORM\_INT\_SUMM}{1}{16}{0}%
\regfield{ABN\_INT\_SUMM}{1}{15}{0}%
\regfield{EARLY\_RECV\_INT}{1}{14}{0}%
\regfield{FATAL\_BUS\_ERR\_INT}{1}{13}{0}%
\regfield{(reserved)}{2}{11}{00}%
\regfield{EARLY\_TRANS\_INT}{1}{10}{0}%
\regfield{RECV\_WDT\_TO}{1}{9}{0}%
\regfield{RECV\_PROC\_STOP}{1}{8}{0}%
\regfield{RECV\_BUF\_UNAVAIL}{1}{7}{0}%
\regfield{RECV\_INT}{1}{6}{0}%
\regfield{TRANS\_UNDFLOW}{1}{5}{0}%
\regfield{RECV\_OVFLOW}{1}{4}{0}%
\regfield{TRANS\_JABBER\_TO}{1}{3}{0}%
\regfield{TRANS\_BUF\_UNAVAIL}{1}{2}{0}%
\regfield{TRANS\_PROC\_STOP}{1}{1}{0}%
\regfield{TRANS\_INT}{1}{0}{0}%
\reglabel{Reset}\regnewline%
\begin{regdesc}\begin{reglist}
\item [EMAC\_LPI\_INT] The raw interrupt status of EMAC\_LPI\_INT. (RO)
\item [TS\_TRI\_INT] The raw interrupt status of TS\_TRI\_INT. (RO)
\item [EMAC\_PMT\_INT] The raw interrupt status of EMAC\_PMT\_INT. (RO)

\item [ERROR\_BITS] Represents the type of error that caused a Bus Error, for example, error response on the AHB interface.\\
0: Error during RX DMA Write Data Transfer\\
1: Error during TX DMA Read Data Transfer\\
2: Error during RX DMA Descriptor Write Access\\
3: Error during TX DMA Descriptor Write Access\\
4: Error during RX DMA Descriptor Read Access\\
5: Error during TX DMA Descriptor Read Access\\
Valid only when Bit 13 (FBI) is 1. This field does not generate an interrupt. (RO)
\item [TRANS\_PROC\_STATE] Represents the Transmit DMA FSM state.\\
0: Stopped; Reset or Stop Transmit Command issued\\
1: Running; Fetching Transmit Transfer Descriptor\\
2: Running; Waiting for status\\
3: Running; Reading Data from Host memory buffer and queuing it to transmit buffer (TX FIFO)\\
4: TIME\_STAMP write state\\
5: Reserved for future use\\
6: Suspended; Transmit Descriptor Unavailable or Transmit Buffer Underflow\\
7: Running; Closing Transmit Descriptor\\
This field does not generate an interrupt. (RO)
\item [Continued on the next page...]
\end{reglist}\end{regdesc}
\end{register}

\addtocounter{Regfloat}{-1}
\begin{register}{H}{DMASTATUS\_REG}{0x1014}
\begin{regdesc}\begin{reglist}
\item [Continued from the previous page...]

\item [NORM\_INT\_SUMM] Represents the logical OR of the following when the corresponding interrupt bits are enabled in Interrupt Enable Register.\\
Bit[0]: Transmit Interrupt\\
Bit[2]: Transmit Buffer Unavailable\\
Bit[6]: Receive Interrupt\\
Bit[14]: Early Receive Interrupt\\
Only unmasked interrupts affect this bit.\\This is a sticky bit and must be cleared (by writing 1 to this bit) each time a corresponding bit, which causes this set to be set, is cleared. (R/SS/WC)
\item [ABN\_INT\_SUMM] Represents the logical OR of the following when the corresponding interrupt bits are enabled in Interrupt Enable Register.\\
Bit[1]: Transmit Process Stopped\\
Bit[3]: Transmit Jabber Timeout\\
Bit[4]: Receive FIFO Overflow\\
Bit[5]: Transmit Underflow\\
Bit[7]: Receive Buffer Unavailable\\
Bit[8]: Receive Process Stopped\\
Bit[9]: Receive Watchdog Timeout\\
Bit[10]: Early Transmit Interrupt\\
Bit[13]: Fatal Bus Error\\
Only unmasked interrupts affect this bit.\\
This is a sticky bit and must be cleared each time a corresponding bit, which causes this bit to be set, is cleared. (R/SS/WC)
\item [EARLY\_RECV\_INT] The raw interrupt status of EARLY\_RECV\_INT.(R/SS/WC)
\item [FATAL\_BUS\_ERR\_INT] The raw interrupt status of FATAL\_BUS\_ERR\_INT. (R/SS/WC)
\item [EARLY\_TRANS\_INT] The raw interrupt status of  EARLY\_TRANS\_INT. (R/SS/WC)
\item [RECV\_WDT\_TO] Represents whether the Receive Watchdog Timer expired while receiving the current frame and the current frame is truncated after the watchdog timeout.\\
0: Not expired\\
1: Expired\\(R/SS/WC)
\item [RECV\_PROC\_STOP] Represents whether the Receive Process enters the Stopped state.\\
0: Not stopped\\
1: Stopped\\(R/SS/WC)

\item [Continued on the next page...]
\end{reglist}\end{regdesc}
\end{register}

\addtocounter{Regfloat}{-1}
\begin{register}{H}{DMASTATUS\_REG}{0x1014}
\begin{regdesc}\begin{reglist}\item [Continued from the previous page...]


\item [RECV\_BUF\_UNAVAIL] Represents whether the Receive Buffer is available.\\
0: Available\\
1: Unavailable\\
This bit is set only when the previous Receive Descriptor is owned by the DMA.\\
When set, the Host owns the Next Descriptor in the Receive List and the DMA cannot acquire it. The Receive Process is suspended. \\
To resume processing Receive descriptors, the Host should change the ownership of the descriptor and issue a Receive Poll Demand command. If no Receive Poll Demand is issued, the Receive Process resumes when the next recognized incoming frame is received. (R/SS/WC)
\item [RECV\_INT] The raw interrupt status of RECV\_INT. (R/SS/WC)
\item [TRANS\_UNDFLOW] Represents whether the Transmit Buffer had an Underflow during frame transmission.\\
0: No underflow\\
1: Underflow\\
When this bit is 1, transmission is suspended and an Underflow Error TDES0[1] is set. (R/SS/WC)
\item [RECV\_OVFLOW] Represents whether the Receive Buffer had an Overflow during frame reception.\\
0: No overflow\\
1: Overflow\\
If the partial frame is transferred to the application, the overflow status is set in RDES0[11]. (R/SS/WC)
\item [TRANS\_JABBER\_TO] Represents whether the Transmit Jabber Timer expired, which happens when the frame size exceeds 2,048 (10,240 bytes when the Jumbo frame is enabled).\\
0: Not expired\\
1: Expired\\
When the Jabber Timeout occurs, the transmission process is aborted and placed in the Stopped state. This causes the Transmit Jabber Timeout TDES0[14] flag to assert. (R/SS/WC)
\item [TRANS\_BUF\_UNAVAIL]  Represents whether the Transmit Buffer is available.\\
0: Available\\
1: Unavailable\\
When this bit is 1, the Host owns the Next Descriptor in the Transmit List and the DMA cannot acquire it. Transmission is suspended. Bits[22:20] explain the Transmit Process state transitions. \\
To resume processing Transmit descriptors, the Host should change the ownership of the descriptor by setting TDES0[31] and then issue a Transmit Poll Demand command. (R/SS/WC)
\item [TRANS\_PROC\_STOP] Represents whether the transmission is stopped.\\
0: Not stopped\\
1: Stopped\\(R/SS/WC)
\item [TRANS\_INT] The raw interrupt status of TRANS\_INT. (R/SS/WC)
\end{reglist}\end{regdesc}
\end{register}


\begin{register}{H}{DMAOPERATION\_MODE\_REG}{0x1018}\label{regdesc:DMAOPERATIONMODEREG}
\regfield{(reserved)}{5}{27}{00000}%
\regfield{DIS\_DROP\_TCPIP\_ERR\_FRAM}{1}{26}{0}%
\regfield{(reserved)}{1}{25}{0}%
\regfield{DIS\_FLUSH\_RECV\_FRAMES}{1}{24}{0}%
\regfield{(reserved)}{3}{21}{000}%
\regfield{FLUSH\_TX\_FIFO}{1}{20}{0}%
\regfield{(reserved)}{3}{17}{0}%
\regfield{TX\_THRESH\_CTRL}{3}{14}{0}%
\regfield{START\_STOP\_TRANSMISSION\_COMMAND}{1}{13}{0}%
\regfield{(reserved)}{2}{11}{0}%
\regfield{(reserved)}{2}{9}{0}%
\regfield{(reserved)}{1}{8}{0}%
\regfield{FWD\_ERR\_FRAME}{1}{7}{0}%
\regfield{FWD\_UNDER\_GF}{1}{6}{0}%
\regfield{DROP\_GFRM}{1}{5}{0}%
\regfield{RX\_THRESH\_CTRL}{2}{3}{0}%
\regfield{OPT\_SECOND\_FRAME}{1}{2}{0}%
\regfield{START\_STOP\_RX}{1}{1}{0}%
\regfield{(reserved)}{1}{0}{0}%
\reglabel{Reset}\regnewline%
\begin{regdesc}\begin{reglist}
\item [DIS\_DROP\_TCPIP\_ERR\_FRAM] Configures whether to disable dropping of TCP/IP checksum error frames.\\
0: Enable\\
1: Disable\\
When FWD\_ERR\_FRAME is 0, this bit is ignored and all error frames are dropped. (R/W)
\item [DIS\_FLUSH\_RECV\_FRAMES] Configures whether to disable flushing of received frames because of the unavailability of receive descriptors or buffers.\\
0: Enable\\
1: Disable\\(R/W)
\item [FLUSH\_TX\_FIFO] Configures whether to flush TX FIFO to default values.\\
0: Not flush\\
1: Flush\\  (R/WS/SC)

\item [TX\_THRESH\_CTRL] Configures the threshold of TX FIFO. Transmission starts when the frame size within the TX FIFO is larger than the threshold. In addition, full frames with a length less than the threshold are also transmitted.\\
0: 64\\
1: 128\\
2: 192\\
3: 256\\
4: 40\\
5: 32\\
6: 24\\
7: 16\\
Transmission starts when the frame size within the TX FIFO is larger than the threshold. In addition, full frames with a length less than the threshold are also transmitted.\\(R/W)

\item[Continued on the next page...]
\end{reglist}\end{regdesc}
\end{register}
\addtocounter{Regfloat}{-1}
\begin{register}{H}{DMAOPERATION\_MODE\_REG}{0x1018}
\begin{regdesc}\begin{reglist}
\item[Continued from the previous page...]

\item [START\_STOP\_TRANSMISSION\_COMMAND] Configures whether to start or stop the transmission.\\
0: Stop\\
1: Start\\ (R/W)
\item [FWD\_ERR\_FRAME] Configures whether RX FIFO forward frames with error status (CRC error, collision error, giant frame, watchdog timeout, or overflow)\\
0: Drop\\
1: Forward\\
(R/W)
\item [FWD\_UNDER\_GF] Configures whether the RX FIFO forward undersized (frames with no Error and length less than 64 bytes) good frames including PAD and CRC.\\
0: Drop\\
1: Forward\\
(R/W)
\item [DROP\_GFRM] Configures whether to the ETH\_MAC drops the received giant frames in the RX FIFO.\\
0: Not drop\\
1: Drop\\(R/W)
\item [RX\_THRESH\_CTRL] Configures the threshold of the RX FIFO.\\
0: 64\\
1: 32\\
2: 96\\
3: 128\\
Transfer (request) to DMA starts when the frame size within the RX FIFO is larger than the threshold.
(R/W)
\item [OPT\_SECOND\_FRAME] Configures whether the DMA processes the second frame of the Transmit data even before the status for the first frame is obtained.\\
0: Not process\\
1: Process\\ (R/W)
\item [START\_STOP\_RX] Configures whether to start or stop RX DMA reception.\\
0: Stop reception after the transfer of the current frame\\
1: Start reception\\(R/W)
\end{reglist}\end{regdesc}
\end{register}

\begin{register}{H}{DMAIN\_EN\_REG}{0x101C}\label{regdesc:DMAINENREG}
\regfield{(reserved)}{15}{17}{000000000000000}%
\regfield{DMAIN\_NISE}{1}{16}{0}%
\regfield{DMAIN\_AISE}{1}{15}{0}%
\regfield{DMAIN\_ERIE}{1}{14}{0}%
\regfield{DMAIN\_FBEE}{1}{13}{0}%
\regfield{(reserved)}{2}{11}{00}%
\regfield{DMAIN\_ETIE}{1}{10}{0}%
\regfield{DMAIN\_RWTE}{1}{9}{0}%
\regfield{DMAIN\_RSE}{1}{8}{0}%
\regfield{DMAIN\_RBUE}{1}{7}{0}%
\regfield{DMAIN\_RIE}{1}{6}{0}%
\regfield{DMAIN\_UIE}{1}{5}{0}%
\regfield{DMAIN\_OIE}{1}{4}{0}%
\regfield{DMAIN\_TJTE}{1}{3}{0}%
\regfield{DMAIN\_TBUE}{1}{2}{0}%
\regfield{DMAIN\_TSE}{1}{1}{0}%
\regfield{DMAIN\_TIE}{1}{0}{0}%
\reglabel{Reset}\regnewline%
\begin{regdesc}\begin{reglist}
\item [DMAIN\_NISE] Write 1 to enable and write 0 to disable the summary of the following normal interrupt enable bits:\\
Bit[0]: Transmit Interrupt\\
Bit[2]: Transmit Buffer Unavailable\\
Bit[6]: Receive Interrupt\\
Bit[14]: Early Receive Interrupt\\
(R/W)
\item [DMAIN\_AISE] Write 1 to enable and write 0 to disable the summary of the following abnormal interrupt enable bis:\\
Bit[1]: Transmit Process Stopped\\
Bit[3]: Transmit Jabber Timeout\\
Bit[4]: Receive FIFO Overflow\\
Bit[5]: Transmit Underflow\\
Bit[7]: Receive Buffer Unavailable\\
Bit[8]: Receive Process Stopped\\
Bit[9]: Receive Watchdog Timeout\\
Bit[10]: Early Transmit Interrupt\\
Bit[13]: Fatal Bus Error\\(R/W)
\item [DMAIN\_ERIE] Write 1 to enable and write 0 to disable EARLY\_RECV\_INT. (R/W)
\item [DMAIN\_FBEE] Write 1 to enable and write 0 to disable FATAL\_BUS\_ERR\_INT.(R/W)
\item [DMAIN\_ETIE] Write 1 to enable and write 0 to disable EARLY\_TRANS\_INT. (R/W)
\item [DMAIN\_RWTE] Write 1 to enable and write 0 to disable RECV\_WDT\_TO\_INT. (R/W)
\item [DMAIN\_RSE] Write 1 to enable and write 0 to disable RECV\_PROC\_STOP\_INT. (R/W)
\item [DMAIN\_RBUE] Write 1 to enable and write 0 to disable RECV\_BUF\_UNAVAIL\_INT. (R/W)
\item [DMAIN\_RIE] Write 1 to enable and write 0 to disable RECV\_INT. (R/W)


\item[Continued on the next page...]
\end{reglist}\end{regdesc}
\end{register}
\addtocounter{Regfloat}{-1}
\begin{register}{H}{DMAIN\_EN\_REG}{0x101C}
\begin{regdesc}\begin{reglist}

\item[Continued from the previous page...]
\item [DMAIN\_UIE] Write 1 to enable and write 0 to disable TRANS\_UNDFLOW\_INT. (R/W)
\item [DMAIN\_OIE] Write 1 to enable and write 0 to disable RECV\_OVFLOW\_INT. (R/W)
\item [DMAIN\_TJTE] Write 1 to enable and write 0 to disable TRANS\_JABBER\_TO\_INT. (R/W)
\item [DMAIN\_TBUE] Write 1 to enable and write 0 to disable TRANS\_BUF\_UNAVAIL\_INT. (R/W)
\item [DMAIN\_TSE] Write 1 to enable and write 0 to disable TRANS\_PROC\_STOP\_INT.(R/W)
\item [DMAIN\_TIE] Write 1 to enable and write 0 to disable TRANS\_INT.(R/W)
\end{reglist}\end{regdesc}
\end{register}


\begin{register}{H}{DMAMISSEDFR\_REG}{0x1020}\label{regdesc:DMAMISSEDFRREG}
\regfield{(reserved)}{3}{29}{00}%
\regfield{Overflow\_BFOC}{1}{28}{{0x0}}%
\regfield{Overflow\_FC}{11}{17}{{0x0}}%
\regfield{Overflow\_BMFC}{1}{16}{{0x0}}%
\regfield{Missed\_FC}{16}{0}{{0x0}}%
\reglabel{Reset}\regnewline%
\begin{regdesc}\begin{reglist}
\item [Overflow\_BFOC] Represents the status of the overflow frame counter.\\
0: No overflow.\\
1: The overflow frame counter (Overflow\_FC) overflows, that is, the RX FIFO overflows with the overflow frame counter at maximum value. In such a scenario, the overflow frame counter is reset to all-zeros.\\ (R/SS/RC)
\item [Overflow\_FC] Represents the number of frames missed by the application.\\
This counter is incremented each time the MTL FIFO overflows. (R/SS/RC)
\item [Overflow\_BMFC] Represents the status of the missed frame counter.\\
0: No overflow.\\
1: The missed frame counter (Missed\_FC) overflows, that is, the Host Receive Buffer being unavailable with the missed frame counter at maximum value. In such a scenario, the Missed frame counter is reset to all-zeros.\\
(R/SS/RC)
\item [Missed\_FC] Represents the number of frames missed by the controller because of the Host Receive Buffer being unavailable.\\
This counter is incremented each time the DMA discards an incoming frame. (R/SS/RC)
\end{reglist}\end{regdesc}
\end{register}

\begin{register}{H}{DMARINTWDTIMER\_REG}{0x1024}\label{regdesc:DMARINTWDTIMERREG}
\regfield{(reserved)}{24}{8}{000000000000000000000000}%
\regfield{RIWTC}{8}{0}{{0x000}}%
\reglabel{Reset}\regnewline%
\begin{regdesc}\begin{reglist}
\item [RIWTC] Configures the number of system clock cycles multiplied by 256.\\
The watchdog timer gets triggered with the programmed value after the RX DMA completes the transfer of a frame for which the RI (RECV\_INT) status bit is not set because of the setting in the corresponding descriptor RDES1[31].\\
When the watchdog timer runs out, the RI bit is set and the timer is stopped. The watchdog timer is reset when the RI bit is set high because of the automatic setting of RI as per RDES1[31] of any received frame. (R/W)
\end{reglist}\end{regdesc}
\end{register}

\begin{register}{H}{DMAAHBSTATUS\_REG}{0x102C}\label{regdesc:DMAHBSTATUSREG}
\regfield{(reserved)}{31}{1}{0000000000000000000000000000000}%
\regfield{AHB\_ST}{1}{0}{{0x0}}%
\reglabel{Reset}\regnewline%
\begin{regdesc}\begin{reglist}
\item [AHB\_ST] Configures the state of the AHB master interface.\\
0: Idle\\
1: Active\\(R/W)
\end{reglist}\end{regdesc}
\end{register}


\begin{register}{H}{DMATXCURRDESC\_REG}{0x1048}\label{regdesc:DMATXCURRDESCREG}
\regfield{}{32}{0}{{0x000000000}}%
\reglabel{Reset}\regnewline%
\begin{regdesc}\begin{reglist}
\item [TRANS\_DECR\_ADDR\_PTR] Represents the address of the current transmit descriptor list updated by the DMA during operation.\\
This field is cleared on reset. (RO)
\end{reglist}\end{regdesc}
\end{register}


\begin{register}{H}{DMARXCURRDESC\_REG}{0x104C}\label{regdesc:DMARXCURRDESCREG}
\regfield{}{32}{0}{{0x000000000}}%
\reglabel{Reset}\regnewline%
\begin{regdesc}\begin{reglist}
\item [RECV\_DECR\_ADDR\_PTR] Represents the address of the current receive descriptor list updated by the DMA during operation.\\
This field is cleared on reset. (RO)
\end{reglist}\end{regdesc}
\end{register}


\begin{register}{H}{DMATXCURRADDR\_BUF\_REG}{0x1050}\label{regdesc:DMATXCURRADDRBUFREG}
\regfield{}{32}{0}{{0x000000000}}%
\reglabel{Reset}\regnewline%
\begin{regdesc}\begin{reglist}
\item [TRANS\_BUFF\_ADDR\_PTR] Represents the address of the Transmit Buffer updated by the DMA during operation.\\
This field is cleared on reset. (RO)
\end{reglist}\end{regdesc}
\end{register}

\begin{register}{H}{DMARXCURRADDR\_BUF\_REG}{0x1054}\label{regdesc:DMARXCURRADDRBUFREG}
\regfield{}{32}{0}{{0x000000000}}%
\reglabel{Reset}\regnewline%
\begin{regdesc}\begin{reglist}
\item [RECV\_BUFF\_ADDR\_PTR] Represents the address of the Receive Buffer updated by the DMA during operation.\\
This field is cleared on reset. (RO)
\end{reglist}\end{regdesc}
\end{register}



\begin{register}{H}{EMACCONFIG\_REG}{0x0000}\label{regdesc:EMACCONFIGREG}
\regfield{(reserved)}{1}{31}{0}%
\regfield{SAIRC}{3}{28}{{0x0}}%
\regfield{ASS2KP}{1}{27}{0}%
\regfield{(reserved)}{1}{26}{0}%
\regfield{EMACCRC\_STRIP}{1}{25}{0}%
\regfield{(reserved)}{1}{24}{0}%
\regfield{EMACWATCHDOG}{1}{23}{0}%
\regfield{EMACJABBER}{1}{22}{0}%
\regfield{(reserved)}{1}{21}{0}%
\regfield{EMACJUMBOFRAME}{1}{20}{0}%
\regfield{EMACINTERFRAMEGAP}{3}{17}{0}%
\regfield{EMACDISABLECRS}{1}{16}{0}%
\regfield{EMACMII}{1}{15}{1}%
\regfield{EMACFESPEED}{1}{14}{0}%
\regfield{EMACRXOWN}{1}{13}{0}%
\regfield{EMACLOOPBACK}{1}{12}{0}%
\regfield{EMACDUPLEX}{1}{11}{0}%
\regfield{EMACRXIPCOFFLOAD}{1}{10}{0}%
\regfield{EMACRETRY}{1}{9}{0}%
\regfield{(reserved)}{1}{8}{0}%
\regfield{EMACPADCRCSTRIP}{1}{7}{0}%
\regfield{EMACBACKOFFLIMIT}{2}{5}{{0x0}}%
\regfield{EMACDEFERRALCHECK}{1}{4}{0}%
\regfield{EMACTX}{1}{3}{0}%
\regfield{EMACRX}{1}{2}{0}%
\regfield{PLTF}{2}{0}{{0x0}}%
\reglabel{Reset}\regnewline%
\begin{regdesc}\begin{reglist}
\item [SAIRC] Configures the source address insertion or replacement for all transmitted frames.\\
Bit 30 specifies which MAC Address register (0 or 1) is used for source address insertion or replacement based on the values of Bits [29:28].\\
Bits [29:28]:\\
0, 1: The input mti\_sa\_ctrl\_i and ati\_sa\_ctrl\_i control the SA field generation\\
2: Insert the content of the MAC Address register (0 or 1) in the SA field of all transmitted frames\\
3: Replace the the content of MAC Address register (0 or 1) in the SA field of all transmitted frames\\
Bit 30:\\
0: Use MAC Address 0 register\\
1: Use MAC Address 1 register\\
(R/W)
\item [ASS2KP] Configures whether the MAC considers received frames of more than 2000 bytes as normal packets or giant frames.\\
0: Giant frames. In this case, if Bit[20] (JE) is 0, the MAC considers all received frames of size more than 1,518 bytes (1,522 bytes for tagged) as Giant frames.\\
1: Normal packets.\\
(R/W)
\item [EMACCRC\_STRIP] Configures whether to strip and drop the last 4 bytes (FCS) of all frames of Ether type (Length/Type field greater than or equal to 1,536) before forwarding the frame to the application.\\
0: Not strip and drop\\
1: Strip and drop\\(R/W)
\item [EMACWATCHDOG] Configures whether to disable the watchdog timer.\\
0: Enable. In this case, the MAC does not allow a receive frame with more than 2,048 bytes (10,240 if JE is set high) or the value programmed in Watchdog Timeout Register. The MAC cuts off any bytes received after the watchdog limit number of bytes.\\
1: Disable. In this case, the MAC can receive frames of up to 16,384 bytes. \\(R/W)
\item [EMACJABBER] Configures whether to disable the jabber timer on the transmitter.\\
0: Enable. In this case, the MAC cuts off the transmitter if the application sends out more than 2,048 bytes of data (10,240 if JE is set high) during transmission.\\
1: Disable. In this case, the MAC can transfer frames of up to 16,384 bytes.\\ (R/W)

\item[Continued on the next page...]
\end{reglist}\end{regdesc}
\end{register}

\addtocounter{Regfloat}{-1}
\begin{register}{H}{EMACCONFIG\_REG}{0x0000}
\begin{regdesc}\begin{reglist}
\item[Continued from the previous page...]

\item [EMACJUMBOFRAME] Configures whether the MAC allows Jumbo frames of 9,018 bytes (9,022 bytes for VLAN tagged frames) without reporting a giant frame error in the receive frame status.\\
0: Not allow\\
1: Allow\\(R/W)
\item [EMACINTERFRAMEGAP] Configures the minimum interframe gap (IFG) between frames during transmission.\\
0: 96 bit times\\
1: 88 bit times\\
2: 80 bit times\\
3: 72 bit times\\
4: 64 bit times\\
5: 56 bit times\\
6: 48 bit times\\
7: 40 bit times\\
In the half-duplex mode, the minimum IFG can be configured only for 64 bit times (IFG = 100). (R/W)
\item [EMACDISABLECRS] Configures whether to disable Carrier Sense during transmission.\\
0: The MAC transmitter generates Loss of Carrier or No Carrier errors because of Carrier Sense and can even abort the transmissions. \\
1: The MAC transmitter ignores the MII CRS signal during frame transmission in the half-duplex mode. This request results in no errors generated because of Loss of Carrier or No Carrier during such transmission.\\
(R/W)
\item [EMACMII] Configures the Ethernet line speed.\\
0: For 1000 Mbps operations\\
1: For 10 or 100 Mbps operations\\
In 10 or 100 Mbps operations, the exact line speed is determined by both this bit and the FES (EMACFESPEED) bit. (R/W)
\item [EMACFESPEED] Configures the speed in the MII and RMII interface.\\
0: 10 Mbps\\
1: 100 Mbps\\(R/W)
\item [EMACRXOWN] Configures whether to disable receive own.\\
0: The MAC receives all packets that are given by the PHY while transmitting.\\
1: the MAC disables the reception of frames when TX\_EN triggers an interrupt in the half-duplex mode.\\
This bit is not applicable if the MAC is operating in the full-duplex mode. (R/W)

\item[Continued on the next page...]
\end{reglist}\end{regdesc}
\end{register}

\addtocounter{Regfloat}{-1}
\begin{register}{H}{EMACCONFIG\_REG}{0x0000}
\begin{regdesc}\begin{reglist}
\item[Continued from the previous page...]

\item [EMACLOOPBACK] Configures whether to enable the loopback mode.\\
0: Disable.\\
1: Enable. In this case, the MII Receive clock input (CLK\_RX) is required for the loopback to work properly, because the Transmit clock is not looped-back internally.\\(R/W)
\item [EMACDUPLEX] Configures whether to enable full-duplex mode, so that the MAC can transmit and receive simultaneously.\\
0: Disable\\
1: Enable\\
This bit is RO with a default value of 1 in the full-duplex-only configuration. (R/W)
\item [EMACRXIPCOFFLOAD] Configures whether to enable the checksum offload function\\
0: Disable\\
1: Enable\\
When the checksum offload function is enabled, the MAC calculates the 16-bit one's complement of the one's complement sum of all received Ethernet frame payloads. It also checks whether the IPv4 Header checksum (assumed to be bytes 25/26 or 29/30 (VLAN-tagged) of the received Ethernet frame) is correct for the received frame and gives the status in the receive status word. The MAC also appends the 16-bit checksum calculated for the IP header datagram payload (bytes after the IPv4 header) and appends it to the Ethernet frame transferred to the application (when Type 2 COE is deselected). (R/W)
\item [EMACRETRY] Configures whether to disable retry, so that the MAC attempts only one transmission, and when a collision occurs on the MII interface the MAC ignores the current frame transmission and reports a Frame Abort with excessive collision error in the transmit frame status.\\
0: Enable\\
1: Disable\\
Valid only in the half-duplex mode. (R/W)
\item [EMACPADCRCSTRIP]
Configures whether to enable automatic Pad or CRC stripping.\\
0: Disable. In this case, the MAC passes all incoming frames, without modifying them, to the Host.\\
1: Enable. In this case, the MAC strips the Pad or FCS field on the incoming frames only if the value of the length field is less than 1,536 bytes. All received frames with length field greater than or equal to 1,536 bytes are passed to the application without stripping the Pad or FCS field.\\ (R/W)

\item[Continued on the next page...]
\end{reglist}\end{regdesc}
\end{register}

\addtocounter{Regfloat}{-1}
\begin{register}{H}{EMACCONFIG\_REG}{0x0000}
\begin{regdesc}\begin{reglist}
\item[Continued from the previous page...]

\item [EMACBACKOFFLIMIT]
Configures the back-off limit determining the random integer number (r) of slot time delays (4,096 bit times for 1000 Mbps and 512 bit times for 10/100 Mbps) for which the MAC waits before rescheduling a transmission attempt during retries after a collision.\\
00: k = min (n, 10)\\
01: k = min (n, 8)\\
10: k = min (n, 4)\\
11: k = min (n, 1)\\
where n = retransmission attempt.\\
The random integer r takes the value in the
range 0 <= r < kth power of 2\\
Valid only in the half-duplex mode. (R/W)
\item [EMACDEFERRALCHECK] Configures whether to enable the deferral check function.\\
0: Disable\\
1: Enable\\(R/W)
\item [EMACTX] Configures whether to enable the transmit state machine of the MAC.\\
0: Disabled after the completion of the transmission of the current frame, and does not transmit any further frames\\
1: Enabled for transmission on the MII\\ (R/W)
\item [EMACRX] Configures whether to enable the receiver state machine of the MAC.
0: Disabled after the completion of the reception of the current frame, and does not receive any further frames from the MII\\
1: Enabled for receiving frames from the MII\\ (R/W)
\item [PLTF] Configures the number of preamble bytes that are added to the beginning of every Transmit frame.\\
0: 7 bytes of preamble\\
1: 5 byte of preamble\\
2: 3 bytes of preamble\\
3: Reserved\\
The preamble reduction occurs only when the MAC is operating in the full-duplex mode. (R/W)
\end{reglist}\end{regdesc}
\end{register}

\begin{register}{H}{EMACFF\_REG}{0x0004}\label{regdesc:EMACFFREG}
\regfield{RECEIVE\_ALL}{1}{31}{0}%
\regfield{(reserved)}{14}{17}{00000000000000}%
\regfield{VTFE}{1}{16}{1}%
\regfield{(reserved)}{6}{10}{000000}%
\regfield{SAFE}{1}{9}{0}%
\regfield{SAIF}{1}{8}{0}%
\regfield{PCF}{2}{6}{{0x0}}%
\regfield{DBF}{1}{5}{0}%
\regfield{PAM}{1}{4}{0}%
\regfield{DAIF}{1}{3}{0}%
\regfield{(reserved)}{2}{1}{00}%
\regfield{PMODE}{1}{0}{0}%
\reglabel{Reset}\regnewline%
\begin{regdesc}\begin{reglist}
\item [RECEIVE\_ALL] Configures whether to pass all received modules.\\
0: The MAC Receiver module passes only those frames to the Application that pass the SA or DA address filter.\\
1: The MAC Receiver module passes all received frames, irrespective of whether they pass the address filter or not, to the Application. The result of the SA or DA filtering is updated (pass or fail) in the corresponding bits in the Receive Status Word.\\(R/W)
\item [VTFE] Configures whether to enable the VLAN Tag filter. \\
0: Disable. In this case, the MAC forwards all frames irrespective of the match status of the VLAN Tag.\\
1: Enable. In this case, the MAC drops VLAN-tagged frames that do not match the VLAN Tag comparison.\\
(R/W)
\item [SAFE] Configures whether to enable the source address filter.\\
0: Disable. In this case, the MAC forwards the received frame to the application with an updated SAF bit of the RX Status depending on the SA address comparison.\\
1: Enable. In this case, the MAC compares the SA field of the received frames with the values programmed in the enabled SA registers. If the comparison fails, the MAC drops the frame.\\(R/W)
\item [SAIF] Configures whether to enable SA inverse filtering.\\
0: Disable. In this case, frames whose SA does not match the SA registers are marked as failing the SA Address filter.\\
1: Enable. In this case, the Address Check block operates in inverse filtering mode for the SA address comparison. The frames whose SA matches the SA registers are marked as failing the SA Address filter.\\(R/W)

\item[Continued on the next page...]
\end{reglist}\end{regdesc}
\end{register}

\addtocounter{Regfloat}{-1}
\begin{register}{H}{EMACFF\_REG}{0x0004}
\begin{regdesc}\begin{reglist}
\item[Continued from the previous page...]

\item [PCF] Configures the forwarding of all control frames (including unicast and multicast PAUSE frames).\\
0: The MAC filters all control frames.\\
1: The MAC forwards all control frames except PAUSE control frames to application even if they fail the Address filter.\\
2: The MAC forwards all control frames to application even if they fail the Address Filter.\\
3: The MAC forwards control frames that pass the Address Filter.\\
The following conditions should be true for the PAUSE control frames processing:\\
Condition 1: The MAC is in the full-duplex mode and flow control is enabled by setting Bit 2 (RFE) of Register 6 (Flow Control Register) to 1.\\
Condition 2: The destination address (DA) of the received frame matches the special multicast address or EMACADDR0 when Bit 3 (UP) of the Register 6 (Flow Control Register) is set.\\
Condition 3: The Type field of the received frame is 0x8808 and the OPCODE field is 0x0001.\\
(R/W)
\item [DBF] Configures whether to disable broadcast frames.\\
0: Enable. The AFM (Address Filtering Module) module passes all received broadcast frames.\\
1: Disable. The AFM module filters all incoming broadcast frames. In addition, it overrides all other filter settings.\\ (R/W)
\item [PAM] Configures whether to pass all multicast frames.\\
0: Filtering of multicast frame depends on HMC bit\\
1: All received frames with a multicast destination address (first bit in the destination address field is '1') are passed\\
(R/W)
\item [DAIF] Configures whether to enable DA inverse filtering for both unicast and multicast frames.\\
0: Disable\\
1: Enable\\
(R/W)
\item [PMODE] Configures whether to enable the promiscuous mode.\\
0: Disable.\\
1: Enabled. In this case, the Address Filter module passes all incoming frames regardless of its destination or source address. The SA or DA Filter Fails status bits of the Receive Status Word are always cleared when PR (PRI\_RATIO) is set.\\(R/W)
\end{reglist}\end{regdesc}
\end{register}

\begin{register}{H}{EMACMIIADDR\_REG}{0x0010}\label{regdesc:EMACMIIADDRREG}
\regfield{(reserved)}{16}{16}{0000000000000000}%
\regfield{MIIDEV}{5}{11}{{0x00}}%
\regfield{MIIREG}{5}{6}{{0x00}}%
\regfield{MIICSRCLK}{4}{2}{{0x00}}%
\regfield{MIIWRITE}{1}{1}{0}%
\regfield{MIIBUSY}{1}{0}{0}%
\reglabel{Reset}\regnewline%
\begin{regdesc}\begin{reglist}
\item [MIIDEV] Configures which of the 32 possible PHY devices are being accessed. (R/W)
\item [MIIREG]\label{fielddesc:MIIREG} Configures the desired address register in the selected PHY device.  (R/W)
\item  [MIICSRCLK]\label{fielddesc:MIICSRCLK} Configures the CSR clock frequency.\\
0: APB clock frequency is 80 MHz, and MDC clock frequency is APB\_CLK/42\\
3: APB clock frequency is 40 MHz, and MDC clock frequency is APB\_CLK/26\\
Other values are reserved\\(R/W) 
\item [MIIWRITE] Configures the direction of the operation that uses \hyperref[regdesc:EMACMIIDATAREG]{MII\_DATA}.\\
0: Read operation\\
1: Write operation\\(R/W)
\item [MIIBUSY]\label{fielddesc:MIIBUSY} Configures the state of the MII.\\
0: Idle\\
1: Busy\\
This bit is used in combination with \hyperref[fielddesc:MIIREG]{MIIREG} and \hyperref[regdesc:EMACMIIDATAREG]{MII\_DATA}.\\
This bit should read logic 0 (default) before writing to \hyperref[fielddesc:MIIREG]{MIIREG} and \hyperref[regdesc:EMACMIIDATAREG]{MII\_DATA}.\\
To read or write to \hyperref[fielddesc:MIIREG]{MIIREG} and \hyperref[regdesc:EMACMIIDATAREG]{MII\_DATA}, software (the user) should set this field to 1.\\
\hyperref[regdesc:EMACMIIDATAREG]{MII\_DATA} should be kept valid (data remains unchanged) when it is accessed until this field is cleared by hardware (the MAC).\\
Note that ESP32-P4 MAC does not receive ACK from PHY during a read or write access to \hyperref[fielddesc:MIIREG]{MIIREG} and \hyperref[regdesc:EMACMIIDATAREG]{MII\_DATA}. (R/WS/SC)
\end{reglist}\end{regdesc}
\end{register}


\begin{register}{H}{EMACMIIDATA\_REG}{0x0014}\label{regdesc:EMACMIIDATAREG}
\regfield{(reserved)}{16}{16}{0000000000000000}%
\regfield{MII\_DATA}{16}{0}{{0x00000}}%
\reglabel{Reset}\regnewline%
\begin{regdesc}\begin{reglist}
\item [MII\_DATA] Configures the 16-bit data value read from the PHY after a Management Read operation or the 16-bit data value to be written to the PHY before a Management Write operation.(R/W)
\end{reglist}\end{regdesc}
\end{register}


\begin{register}{H}{EMACFC\_REG}{0x0018}\label{regdesc:EMACFLOWCONTROLREG}
\regfield{PAUSE\_TIME}{16}{16}{{0x00000}}%
\regfield{(reserved)}{8}{8}{00000000}%
\regfield{DZPQ}{1}{7}{{0}}%
\regfield{(reserved)}{1}{6}{0}%
\regfield{PLT}{2}{4}{{0x0}}%
\regfield{UPFD}{1}{3}{0}%
\regfield{RFCE}{1}{2}{0}%
\regfield{TFCE}{1}{1}{0}%
\regfield{FCBBA}{1}{0}{0}%
\reglabel{Reset}\regnewline%
\begin{regdesc}\begin{reglist}
\item [PAUSE\_TIME] Configures the value to be used in the Pause Time field in the transmit control frame. If the Pause Time bits is configured to be double-synchronized to the (G)MII clock domain, then consecutive writes to this register should be performed only after at least four clock cycles in the destination clock domain. (R/W)
\item [DZPQ] Configures whether to disable the automatic generation of the Zero-Quanta Pause Control frames.\\
0: Enable\\
1: Disable\\(R/W)
\item [PLT] Configures the threshold of the PAUSE timer at which the PAUSE Frame is automatically retransmitted.\\
0: The threshold is Pause time minus 4 slot times (PT - 4 slot times).\\
1: The threshold is Pause time minus 28 slot times (PT - 28 slot times).\\
2: The threshold is Pause time minus 144 slot times (PT - 144 slot times).\\
3: The threshold is Pause time minus 256 slot times (PT - 256 slot times).\\
The threshold values should be always less than the Pause Time configured in Bits[31:16]. For example, if PT = 100H (256 slot-times), and PLT = 01, then a second PAUSE frame is automatically transmitted 228 (256 - 28) slot times after the first PAUSE frame is transmitted.\\ (R/W)
\item [UPFD] Configures whether to enable unicast pause frame detection.\\
0: Disable. The MAC only detects Pause frames with a unique multicast address.\\
1: Enable. The MAC can also detect Pause frames with unicast address of the station. This unicast address should be as specified in the EMACADDR0 High Register and EMACADDR0 Low Register.\\(R/W)
\item [RFCE] Configures whether to enable the decode function of the Pause time.\\
0: Disable\\
1: Enable. The MAC decodes the received Pause frame and disables its transmitter for a specified (Pause) time\\(R/W)
\item [TFCE] Configures whether to enable transmit flow control in full-duplex mode, and whether to enable the back-pressure feature in half-duplex mode.\\
0: Disable\\
1: Enable\\
(R/W)

\item[Continued on the next page...]
\end{reglist}\end{regdesc}
\end{register}

\addtocounter{Regfloat}{-1}
\begin{register}{H}{EMACFC\_REG}{0x0018}
\begin{regdesc}\begin{reglist}
\item[Continued from the previous page...]

\item [FCBBA]
Configures whether to enable a Pause Control frame in the full-duplex mode, and whether to activate the back-pressure function in the half-duplex mode if the TFCE bit is set.\\
0: Disable\\
1: Enable\\
In the full-duplex mode, this bit should be read as 0 before writing to the Flow Control register. To initiate a Pause control frame, the Application must set this bit to 1. During a transfer of the Control Frame, this bit continues to be set to signify that a frame transmission is in progress. After the completion of Pause control frame transmission, the MAC resets this bit to 0. The Flow Control register should not be written to until this bit is cleared. (R/WS/SC)\\
In the half-duplex mode, when this bit is set (and TFCE is set), then backpressure is asserted by the MAC. During backpressure, when the MAC receives a new frame, the transmitter starts sending a JAM pattern resulting in a collision. When the MAC is configured for the full-duplex mode, the BPA is automatically disabled. (R/W)
\end{reglist}\end{regdesc}
\end{register}

\begin{register}{H}{EMACVLANTAG\_REG}{0x001C}\label{regdesc:EMACVLANTAGREG}
\regfield{(reserved)}{13}{19}{0000000000000}%
\regfield{ESVL}{1}{18}{0}%
\regfield{VTIM}{1}{17}{{0}}%
\regfield{ETV}{1}{16}{0}%
\regfield{VL}{16}{0}{{0x0000}}%
\reglabel{Reset}\regnewline%
\begin{regdesc}\begin{reglist}
\item [ESVL] Configures whether to enable S-VLAN.\\
0: Disable.\\
1: Enable. The MAC transmitter and receiver also consider the S-VLAN (Type = 0x88A8) frames as valid VLAN tagged frames.\\(R/W)
\item [VTIM] Configures whether to enable VLAN Tag inverse match.\\
0: Disable. The frames with matched VLAN Tag are marked as matched.\\
1: Enable. The frames that do not have matching VLAN Tag are marked as matched.\\ (R/W)
\item [ETV] Configures whether to enable 12-bit VLAN Tag comparison.\\
0: Disable. All 16 bits of the 15th and 16th bytes of the received VLAN frame are used for comparison and VLAN hash filtering.\\
1: Enable. A 12-bit VLAN identifier is used for comparing and filtering instead of the complete 16-bit VLAN tag. Bits [11:0] of VLAN tag are compared with the corresponding field in the received VLAN-tagged frame.\\(R/W)
\item [VL] Configures the 802.1Q VLAN tag to identify the VLAN frames and is compared to the 15th and 16th bytes of the frames being received for VLAN frames.\\
When the ETV bit is set, only the VID (Bits[11:0]) is used for comparison. \\
If VL (VL[11:0] if ETV is set) is all zeros, the MAC does not check the fifteenth and 16th bytes for VLAN tag comparison, and declares all frames with a Type field value of 0x8100 or 0x88a8 as VLAN frames.(R/W)
\end{reglist}\end{regdesc}
\end{register}


\begin{register}{H}{EMACDEBUG\_REG}{0x0024}\label{regdesc:EMACDEBUGREG}
\regfield{(reserved)}{6}{26}{000000}%
\regfield{MTLTSFFS}{1}{25}{0}%
\regfield{MTLTFNES}{1}{24}{0}%
\regfield{(reserved)}{1}{23}{0}%
\regfield{MTLTFWCS}{1}{22}{0}%
\regfield{MTLTFRCS}{2}{20}{{0x0}}%
\regfield{MACTP}{1}{19}{0}%
\regfield{MACTFCS}{2}{17}{{0x0}}%
\regfield{MACTPES}{1}{16}{0}%
\regfield{(reserved)}{6}{10}{000000}%
\regfield{MTLRFFLS}{2}{8}{{0x0}}%
\regfield{(reserved)}{1}{7}{0}%
\regfield{MTLRFRCS}{2}{5}{{0x0}}%
\regfield{MTLRFWCAS}{1}{4}{0}%
\regfield{(reserved)}{1}{3}{0}%
\regfield{MACRFFCS}{2}{1}{{0x0}}%
\regfield{MACRPES}{1}{0}{0}%
\reglabel{Reset}\regnewline%
\begin{regdesc}\begin{reglist}
\item [MTLTSFFS] Represents whether MTL TX Status FIFO is full and can accept any more frames for transmission.\\
0: Not Full\\
1: Full\\ (RO)
\item [MTLTFNES] Represents whether MTL TX FIFO is not empty and some data is left for transmission.\\
0: Empty\\
1: Not empty\\(RO)
\item [MTLTFWCS] Represents whether MTL TX FIFO  Write Controller is active and transferring data to the TX FIFO.\\
0: Inactive\\
1: Active\\(RO)
\item [MTLTFRCS] Represents the state of the TX FIFO Read Controller.\\
0: IDLE state\\
1: READ state (transferring data to MAC transmitter)\\
2: Waiting for TXStatus from MAC transmitter\\
3: Writing the received TXStatus or flushing the TX FIFO\\
(RO)
\item [MACTP] Represents whether the MAC transmitter is in the PAUSE condition (in the full-duplex only mode) and hence does not schedule any frame for transmission.\\
0: Not in pause\\
1: In pause\\
(RO)
\item [MACTFCS] Represents the state of the MAC Transmit Frame Controller module.\\
0: IDLE state\\
1: Waiting for Status of the previous frame or IFG or backoff period to be over\\
2: Generating and transmitting a PAUSE control frame (in the full-duplex mode)\\
3: Transferring input frame for transmission\\
(RO)
\item [MACTPES] Represents the status of the MAC MII transmit protocol engine.\\
0: IDLE state\\
1: Actively transmitting data\\(RO)

\item[Continued on the next page...]
\end{reglist}\end{regdesc}
\end{register}
\addtocounter{Regfloat}{-1}
\begin{register}{H}{EMACDEBUG\_REG}{0x0024}
\begin{regdesc}\begin{reglist}
\item[Continued from the previous page...]

\item [MTLRFFLS] Represents the status of the fill-level of the RX FIFO.\\
0: RX FIFO Empty\\
1: RX FIFO fill level is below the flow-control deactivate threshold\\
2: RX FIFO fill level is above the flow-control activate threshold\\
3: RX FIFO Full\\
(RO)
\item [MTLRFRCS] Represents the state of the RX FIFO read Controller.\\
0: IDLE state\\
1: Reading frame data\\
2: Reserved\\
3: Flushing the frame data and status\\
(RO)
\item [MTLRFWCAS] Represents the state of the MTL RX FIFO Write Controller.\\
0: Inactive\\
1: Active and is transferring a received frame to the FIFO.\\(RO)
\item [MACRFFCS] Represents the state of the small FIFO Read (bit 1) and Write (bit 0) controllers of the MAC Receive Frame Controller Module.\\
0: Inactive\\
1: Active\\(RO)
\item [MACRPES] Represents the state of the MAC GMII or MII receive protocol engine.\\
0: IDLE\\
1: Actively receiving data\\(RO)
\end{reglist}\end{regdesc}
\end{register}


\begin{register}{H}{PMT\_RWUFFR\_REG}{0x0028}\label{regdesc:PMTRWUFFRREG}
\regfield{PMT\_RWUFFR}{32}{0}{00000000000000000000000000000000}%
\reglabel{Reset}\regnewline%
\begin{regdesc}\begin{reglist}
\item [WKUPPKTFILTER] When RWKPTR is 0 \verb+~+ 3, configures the masking of filter 0 \verb+~+ 3 wakeup frame.\\
Bit[30:0] corresponds to offset+j (j = 0 \verb+~+ 30).\\ Bit[31] has to be 0.
0: Unmask\\
1: Mask\\
When RWKPTR is 4, bit[3]/bit[11]/bit[19]/bit[27] configure the target address type, and bit[0]/bit[8]/bit[16]/bit[24] configure whether to enable filter 0 \verb+~+ 3.\\
0: Unicast/Disable\\
1: Multicast/Enable\\
When RWKPTR is 5, configures the offset of the filter 0 \verb+~+ 3 masked byte.\\
Bit[7:0] configures the offset of filter 0 masked byte, and [31:24] configures the offset of filter 3 masked byte.\\
When RWKPTR is 6 \verb+~+ 7, configures the 16-bit CRC value of filter 0 \verb+~+ 3 masked byte starting from offset based on the formula:\\
G(x) = x$^\text{16}$ + x$^\text{15}$ + x$^\text{2}$ + 1\\
Bit[7:0] configures the CRC value of filter 0 or filter 2 masked byte when RWKPTR is 6 or 7 respectively, and bit[31:16] configures the CRC value of filter 1 or filter 3 masked byte when RWKPTR is 6 or 7 respectively.\\
(R/W)
\end{reglist}\end{regdesc}
\end{register}


\begin{register}{H}{PMT\_CSR\_REG}{0x002C}\label{regdesc:PMTCSRREG}
\regfield{RWKFILTRST}{1}{31}{0}%
\regfield{(reserved)}{2}{29}{00}%
\regfield{RWKPTR}{5}{24}{00000}%
\regfield{(reserved)}{14}{10}{00000000000000}%
\regfield{GLBLUCAST}{1}{9}{00}%
\regfield{(reserved)}{2}{7}{00}%
\regfield{RWKPRCVD}{1}{6}{0}%
\regfield{MGKPRCVD}{1}{5}{0}%
\regfield{(reserved)}{2}{3}{00}%
\regfield{RWKPKTEN}{1}{2}{0}%
\regfield{MGKPKTEN}{1}{1}{0}%
\regfield{PWRDWN}{1}{0}{0}%
\reglabel{Reset}\regnewline%
\begin{regdesc}\begin{reglist}
\item [RWKFILTRST] Configures whether to reset the RWKPTR register.\\
0: Release from reset\\
1: Reset\\ (R/WS/SC)
\item [RWKPTR] Configures the purpose of the PMT\_RWUFFR register. For details, see the description of PMT\_RWUFFR. (RO)
\item [GLBLUCAST] Configures whether to enable the unicast frame of target address filter as the wakeup frame.\\
0: Disable\\
1: Enable\\ (R/W)
\item [RWKPRCVD] Represents whether a remote wakeup frame event has been received.\\
0: Not received\\
1: Received\\ (R/SS/RC)
\item [MGKPRCVD]  Represents whether a magic frame event has been received.\\
0: Not received\\
1: Received\\ (R/SS/RC)
\item [RWKPKTEN] Configures whether to enable power management events using remote wakeup frames.\\
0: Disable\\
1: Enable\\ (R/W)
\item [MGKPKTEN] Configures whether to enable power management events using magic frames.\\
0: Disable\\
1: Enable\\(R/W)
\item [PWRDWN] Configures whether the EMAC receiver drops all received frames until it receives the expected magic frame or remote wakeup frame.\\
Valid only when MGKPKTEN, GLBLUCAST or RWKPKTEN is 1. (R/WS/SC)


\end{reglist}\end{regdesc}
\end{register}


\begin{register}{H}{EMACLPI\_CSR\_REG}{0x0030}\label{regdesc:EMACLPICSRREG}
\regfield{(reserved)}{12}{20}{00000000000}%
\regfield{LPITXA}{1}{19}{0}%
\regfield{(reserved)}{1}{18}{0}%
\regfield{PLS}{1}{17}{0}%
\regfield{LPIEN}{1}{16}{0}%
\regfield{(reserved)}{6}{10}{000000}%
\regfield{RLPIST}{1}{9}{0}%
\regfield{TLPIST}{1}{8}{0}%
\regfield{(reserved)}{4}{4}{0}%
\regfield{RLPIEX}{1}{3}{0}%
\regfield{RLPIEN}{1}{2}{0}%
\regfield{TLPIEX}{1}{1}{0}%
\regfield{TLPIEN}{1}{0}{0}%
\reglabel{Reset}\regnewline%
\begin{regdesc}\begin{reglist}
\item [LPITXA] Configures the behavior of the MAC when it is entering or coming out of the LPI mode on the transmit side.\\
0: This bit directly controls behavior of the MAC when it is entering or coming out of the LPI mode.\\
1: This bit controls behavior of the MAC when it is entering or coming out of the LPI mode in combination with LPIEN. When LPIEN is also 1, the MAC enters the LPI mode only after all outstanding frames (in the core) and pending frames (in the application interface) have been transmitted. The MAC comes out of the LPI mode when the application sends any frame for transmission\\(R/W)
\item [PLS] Configures the link status of the PHY.\\
0: Link up\\
1:  Link down\\ (R/W)
\item [LPIEN] Configures whether the EMAC transmitter enters the LPI state.\\
0: Exit\\
1: Enter\\
This bit is cleared when the LPITXA bit is set and the MAC exits the LPI state because of the arrival of a new packet for transmission.\\ (R/W/SC)
\item [RLPIST] Represents the receive LPI state.\\
0: Inactive\\
1: The EMAC is receiving the LPI pattern\\ (RO)
\item [TLPIST] Represents the transmit LPI state.\\
0: Inactive\\
1: The EMAC is transmitting the LPI pattern\\ (RO)
\item [RLPIEX] Represents the receive LPI exit state.\\
0: Inactive\\
1: The EMAC Receiver has stopped receiving the LPI pattern, exited the LPI state, and resumed the normal reception\\ (R/SS/RC)

\item[Continued on the next page...]
\end{reglist}\end{regdesc}
\end{register}
\addtocounter{Regfloat}{-1}
\begin{register}{H}{EMACLPI\_CSR\_REG}{0x0030}
\begin{regdesc}\begin{reglist}
\item[Continued from the previous page...]

\item [RLPIEN] Represents the receive LPI entry state.\\
0: Inactive\\
1: The EMAC Receiver has received an LPI pattern and entered the LPI state\\ (R/SS/RC)
\item [TLPIEX] Represents the transmit LPI exit state.\\
0: Inactive\\
1: The EMAC transmitter has exited the LPI state after the user has cleared the LPIEN bit and the LPI\_TW\_TIMER has expired\\ (R/SS/RC)
\item [TLPIEN] Represents the transmit LPI entry state.\\
0: Inactive\\
1: The EMAC Transmitter has entered the LPI state because of the setting of the LPIEN bit\\(R/SS/RC)
\end{reglist}\end{regdesc}
\end{register}


\begin{register}{H}{EMACLPITIMERSCONTROL\_REG}{0x0034}\label{regdesc:EMACLPITIMERSCONTROLREG}
\regfield{(reserved)}{6}{26}{000000}%
\regfield{LPI\_LS\_TIMER}{10}{16}{0x3E8}%
\regfield{LPI\_TW\_TIMER}{16}{0}{000000000000000}%
\reglabel{Reset}\regnewline%
\begin{regdesc}\begin{reglist}
\item [LPI\_LS\_TIMER] Configures the minimum time for which the link status from the PHY should be up.\\
Measurement unit: millisecond.\\
The default value is 1000 (1 second) as defined in the IEEE standard.\\ (R/W)
\item [LPI\_TW\_TIMER] Configures the minimum time for which the MAC waits after it stops transmitting the LPI pattern to the PHY and before it resumes the normal transmission.\\
Measurement unit: millisecond.\\
The TLPIEX status bit is set after the expiry of this timer.\\ (R/W)
\end{reglist}\end{regdesc}
\end{register}




\begin{register}{H}{EMACINTS\_REG}{0x0038}\label{regdesc:EMACINTSREG}
\regfield{(reserved)}{21}{11}{000000000000000000000}%
\regfield{LPIINTS}{1}{10}{0}%
\regfield{TINTS}{1}{9}{0}%
\regfield{(reserved)}{4}{5}{0000}%
\regfield{PMTINTS}{1}{3}{0}%
\regfield{(reserved)}{3}{0}{000}%
\reglabel{Reset}\regnewline%
\begin{regdesc}\begin{reglist}
\item [LPIINTS] The raw interrupt status of EMAC\_LPI\_INT. (RO)
\item [TINTS] The raw interrupt status of TS\_TRI\_INT. (R/SS/RC)
\item [PMTINTS] The raw interrupt status of EMAC\_PMP\_INT. (RO)
\end{reglist}\end{regdesc}
\end{register}


\begin{register}{H}{EMACINTMASK\_REG}{0x003C}\label{regdesc:EMACINTMASKREG}
\regfield{(reserved)}{21}{11}{000000000000000000000}%
\regfield{LPIINTMASK}{1}{10}{0}%
\regfield{TSINTMASK}{1}{9}{0}%
\regfield{(reserved)}{5}{4}{00000}%
\regfield{PMTINTMASK}{1}{3}{0}%
\regfield{(reserved)}{3}{0}{000}%
\reglabel{Reset}\regnewline%
\begin{regdesc}\begin{reglist}
\item [LPIINTMASK] Write 1 to mask EMAC\_LPI\_INT.(R/W)
\item [TSINTMASK] Write 1 to mask TS\_TRI\_INT. (R/W)
\item [PMTINTMASK] Write 1 to mask EMAC\_PMP\_INT. (R/W)
\end{reglist}\end{regdesc}
\end{register}


\begin{register}{H}{EMACADDR0HIGH\_REG}{0x0040}\label{regdesc:EMACADDR0HIGHREG}
\regfield{ADDRESS\_ENABLE0}{1}{31}{1}%
\regfield{(reserved)}{15}{16}{000000000000000}%
\regfield{MAC\_ADDRESS0\_HI}{16}{0}{{0x0FFFF}}%
\reglabel{Reset}\regnewline%
\begin{regdesc}\begin{reglist}
\item [ADDRESS\_ENABLE0] This bit is always set to 1. (RO)
\item [MAC\_ADDRESS0\_HI] Configures the upper 16 bits of the first 6-byte MAC address [47:32].\\The MAC uses this field for filtering the received frames and inserting the MAC address in the Transmit Flow Control (PAUSE) Frames. (R/W)
\end{reglist}\end{regdesc}
\end{register}


\begin{register}{H}{EMACADDR0LOW\_REG}{0x0044}\label{regdesc:EMACADDR0LOWREG}
\regfield{}{32}{0}{{0x0FFFFFFFF}}%
\reglabel{Reset}\regnewline%
\begin{regdesc}\begin{reglist}
\item [EMACADDR0LOW\_REG] Configures the lower 32 bits of the first 6-byte MAC address.\\ The MAC uses this field for filtering the received frames and inserting the MAC address in the Transmit Flow Control (PAUSE) Frames. (R/W)
\end{reglist}\end{regdesc}
\end{register}


\begin{register}{H}{EMACADDR1HIGH\_REG}{0x0048}\label{regdesc:EMACADDR1HIGHREG}
\regfield{ADDRESS\_ENABLE1}{1}{31}{0}%
\regfield{SOURCE\_ADDRESS}{1}{30}{0}%
\regfield{MASK\_BYTE\_CONTROL}{6}{24}{{0x00}}%
\regfield{(reserved)}{8}{16}{00000000}%
\regfield{MAC\_ADDRESS1\_HI}{16}{0}{{0x0FFFF}}%
\reglabel{Reset}\regnewline%
\begin{regdesc}\begin{reglist}
\item [ADDRESS\_ENABLE1] Configures whether to enable perfect filtering using the second MAC address.\\
0: Disable\\
1: Enable\\ (R/W)
\item [SOURCE\_ADDRESS] Configures whether to which fields of the received frame EMACADDR1 [47:0] compares.\\
0: The DA fields\\
1: The SA fields\\ (R/W)
\item [MASK\_BYTE\_CONTROL] Configures whether to mask MAC address byte bytes when comparing the received DA or SA with the contents of EMACADDR1.\\
0: Unmask\\
1: Mask\\
The correspondence between the mask control bits and the bytes:\\
Bit[29]: EMACADDR1HIGH\_REG[15:8]\\
Bit[28]: EMACADDR1HIGH\_REG[7:0]\\
Bit[27]: EMACADDR1LOW\_REG[31:24]\\
Bit[24]: EMACADDR1LOW\_REG[7:0]\\
You can filter a group of addresses (known as group address filtering) by masking one or more bytes of the address.(R/W)
\item [MAC\_ADDRESS1\_HI] Configures the upper 16 bits of the second 6-byte MAC address [47:32]. (R/W)
\end{reglist}\end{regdesc}
\end{register}


\begin{register}{H}{EMACADDR1LOW\_REG}{0x004C}\label{regdesc:EMACADDR1LOWREG}
\regfield{}{32}{0}{{0x0FFFFFFFF}}%
\reglabel{Reset}\regnewline%
\begin{regdesc}\begin{reglist}
\item [EMACADDR1LOW\_REG] Configures the lower 32 bits of the second 6-byte MAC address.\\The content of this field is undefined until loaded by the Application after the initialization process. (R/W)
\end{reglist}\end{regdesc}
\end{register}


\begin{register}{H}{EMACADDR2HIGH\_REG}{0x0050}\label{regdesc:EMACADDR2HIGHREG}
\regfield{ADDRESS\_ENABLE2}{1}{31}{0}%
\regfield{SOURCE\_ADDRESS2}{1}{30}{0}%
\regfield{MASK\_BYTE\_CONTROL2}{6}{24}{{0x00}}%
\regfield{(reserved)}{8}{16}{00000000}%
\regfield{MAC\_ADDRESS2\_HI}{16}{0}{{0x0FFFF}}%
\reglabel{Reset}\regnewline%
\begin{regdesc}\begin{reglist}
\item [ADDRESS\_ENABLE2] Configures whether to enable perfect filtering using the third MAC address.\\
0: Disable\\
1: Enable\\ (R/W)
\item [SOURCE\_ADDRESS2] Configures whether to which fields of the received frame EMACADDR2 [47:0] compares.\\
0: The DA fields\\
1: The SA fields\\ (R/W)
\item [MASK\_BYTE\_CONTROL2] Configures whether to mask MAC address byte bytes when comparing the received DA or SA with the contents of EMACADDR2.\\
0: Unmask\\
1: Mask\\
The correspondence between the mask control bits and the bytes:\\
Bit[29]: EMACADDR2HIGH\_REG[15:8]\\
Bit[28]: EMACADDR2HIGH\_REG[7:0]\\
Bit[27]: EMACADDR2LOW\_REG[31:24]\\
Bit[24]: EMACADDR2LOW\_REG[7:0]\\
You can filter a group of addresses (known as group address filtering) by masking one or more bytes of the address.(R/W)
\item [MAC\_ADDRESS2\_HI] Configures the upper 16 bits of the third 6-byte MAC address [47:32]. (R/W)
\end{reglist}\end{regdesc}
\end{register}


\begin{register}{H}{EMACADDR2LOW\_REG}{0x0054}\label{regdesc:EMACADDR2LOWREG}
\regfield{}{32}{0}{{0x0FFFFFFFF}}%
\reglabel{Reset}\regnewline%
\begin{regdesc}\begin{reglist}
\item [EMACADDR2LOW\_REG] Configures the lower 32 bits of the third 6-byte MAC address.\\The content of this field is undefined until loaded by the Application after the initialization process.(R/W)
\end{reglist}\end{regdesc}
\end{register}

\begin{register}{H}{EMACADDR3HIGH\_REG}{0x0058}\label{regdesc:EMACADDR3HIGHREG}
\regfield{ADDRESS\_ENABLE3}{1}{31}{0}%
\regfield{SOURCE\_ADDRESS3}{1}{30}{0}%
\regfield{MASK\_BYTE\_CONTROL3}{6}{24}{{0x00}}%
\regfield{(reserved)}{8}{16}{00000000}%
\regfield{MAC\_ADDRESS3\_HI}{16}{0}{{0x0FFFF}}%
\reglabel{Reset}\regnewline%
\begin{regdesc}\begin{reglist}
\item [ADDRESS\_ENABLE3] Configures whether to enable perfect filtering using the fourth MAC address.\\
0: Disable\\
1: Enable\\ (R/W)
\item [SOURCE\_ADDRESS3] Configures whether to which fields of the received frame EMACADDR3 [47:0] compares.\\
0: The DA fields\\
1: The SA fields\\ (R/W)
\item [MASK\_BYTE\_CONTROL3] Configures whether to mask MAC address byte bytes when comparing the received DA or SA with the contents of EMACADDR3.\\
0: Unmask\\
1: Mask\\
The correspondence between the mask control bits and the bytes:\\
Bit[29]: EMACADDR3HIGH\_REG[15:8]\\
Bit[28]: EMACADDR3HIGH\_REG[7:0]\\
Bit[27]: EMACADDR3LOW\_REG[31:24]\\
Bit[24]: EMACADDR3LOW\_REG[7:0]\\
You can filter a group of addresses (known as group address filtering) by masking one or more bytes of the address. (R/W)
\item [MAC\_ADDRESS3\_HI] Configures the upper 16 bits of the fourth 6-byte MAC address [47:32]. (R/W)
\end{reglist}\end{regdesc}
\end{register}


\begin{register}{H}{EMACADDR3LOW\_REG}{0x005C}\label{regdesc:EMACADDR3LOWREG}
\regfield{}{32}{0}{{0x0FFFFFFFF}}%
\reglabel{Reset}\regnewline%
\begin{regdesc}\begin{reglist}
\item [EMACADDR3LOW\_REG] Configures the lower 32 bits of the fourth 6-byte MAC address.\\The content of this field is undefined until loaded by the Application after the initialization process. (R/W)
\end{reglist}\end{regdesc}
\end{register}

\begin{register}{H}{EMACADDR4HIGH\_REG}{0x0060}\label{regdesc:EMACADDR4HIGHREG}
\regfield{ADDRESS\_ENABLE4}{1}{31}{0}%
\regfield{SOURCE\_ADDRESS4}{1}{30}{0}%
\regfield{MASK\_BYTE\_CONTROL4}{6}{24}{{0x00}}%
\regfield{(reserved)}{8}{16}{00000000}%
\regfield{MAC\_ADDRESS4\_HI}{16}{0}{{0x0FFFF}}%
\reglabel{Reset}\regnewline%
\begin{regdesc}\begin{reglist}
\item [ADDRESS\_ENABLE4] Configures whether to enable perfect filtering using the fifth MAC address.\\
0: Disable\\
1: Enable\\ (R/W)
\item [SOURCE\_ADDRESS4] Configures whether to which fields of the received frame EMACADDR4 [47:0] compares.\\
0: The DA fields\\
1: The SA fields\\ (R/W)
\item [MASK\_BYTE\_CONTROL4] Configures whether to mask MAC address byte bytes when comparing the received DA or SA with the contents of EMACADDR4.\\
0: Unmask\\
1: Mask\\
The correspondence between the mask control bits and the bytes:\\
Bit[29]: EMACADDR4HIGH\_REG[15:8]\\
Bit[28]: EMACADDR4HIGH\_REG[7:0]\\
Bit[27]: EMACADDR4LOW\_REG[31:24]\\
Bit[24]: EMACADDR4LOW\_REG[7:0]\\
You can filter a group of addresses (known as group address filtering) by masking one or more bytes of the address.(R/W)
\item [MAC\_ADDRESS4\_HI] Configures the upper 16 bits of the fifth 6-byte MAC address [47:32]. (R/W)
\end{reglist}\end{regdesc}
\end{register}


\begin{register}{H}{EMACADDR4LOW\_REG}{0x0064}\label{regdesc:EMACADDR4LOWREG}
\regfield{}{32}{0}{{0x0FFFFFFFF}}%
\reglabel{Reset}\regnewline%
\begin{regdesc}\begin{reglist}
\item [EMACADDR4LOW\_REG] Configures the lower 32 bits of the fifth 6-byte MAC address.\\The content of this field is undefined until loaded by the Application after the initialization process. (R/W)
\end{reglist}\end{regdesc}
\end{register}

\begin{register}{H}{EMACADDR5HIGH\_REG}{0x0068}\label{regdesc:EMACADDR5HIGHREG}
\regfield{ADDRESS\_ENABLE5}{1}{31}{0}%
\regfield{SOURCE\_ADDRESS5}{1}{30}{0}%
\regfield{MASK\_BYTE\_CONTROL5}{6}{24}{{0x00}}%
\regfield{(reserved)}{8}{16}{00000000}%
\regfield{MAC\_ADDRESS5\_HI}{16}{0}{{0x0FFFF}}%
\reglabel{Reset}\regnewline%
\begin{regdesc}\begin{reglist}
\item [ADDRESS\_ENABLE5] Configures whether to enable perfect filtering using the sixth MAC address.\\
0: Disable\\
1: Enable\\ (R/W)
\item [SOURCE\_ADDRESS5] Configures whether to which fields of the received frame EMACADDR5 [47:0] compares.\\
0: The DA fields\\
1: The SA fields\\ (R/W)
\item [MASK\_BYTE\_CONTROL5] Configures whether to mask MAC address byte bytes when comparing the received DA or SA with the contents of EMACADDR5.\\
0: Unmask\\
1: Mask\\
The correspondence between the mask control bits and the bytes:\\
Bit[29]: EMACADDR5HIGH\_REG[15:8]\\
Bit[28]: EMACADDR5HIGH\_REG[7:0]\\
Bit[27]: EMACADDR5LOW\_REG[31:24]\\
Bit[24]: EMACADDR5LOW\_REG[7:0]\\
You can filter a group of addresses (known as group address filtering) by masking one or more bytes of the address.(R/W)
\item [MAC\_ADDRESS5\_HI] Configures the upper 16 bits of the sixth 6-byte MAC address [47:32]. (R/W)
\end{reglist}\end{regdesc}
\end{register}

\begin{register}{H}{EMACADDR5LOW\_REG}{0x006C}\label{regdesc:EMACADDR5LOWREG}
\regfield{}{32}{0}{{0x0FFFFFFFF}}%
\reglabel{Reset}\regnewline%
\begin{regdesc}\begin{reglist}
\item [EMACADDR5LOW\_REG] Configures the lower 32 bits of the sixth 6-byte MAC address.\\The content of this field is undefined until loaded by the Application after the initialization process. (R/W)
\end{reglist}\end{regdesc}
\end{register}

\begin{register}{H}{EMACADDR6HIGH\_REG}{0x0070}\label{regdesc:EMACADDR6HIGHREG}
\regfield{ADDRESS\_ENABLE6}{1}{31}{0}%
\regfield{SOURCE\_ADDRESS6}{1}{30}{0}%
\regfield{MASK\_BYTE\_CONTROL6}{6}{24}{{0x00}}%
\regfield{(reserved)}{8}{16}{00000000}%
\regfield{MAC\_ADDRESS6\_HI}{16}{0}{{0x0FFFF}}%
\reglabel{Reset}\regnewline%
\begin{regdesc}\begin{reglist}
\item [ADDRESS\_ENABLE6] Configures whether to enable perfect filtering using the seventh MAC address.\\
0: Disable\\
1: Enable\\ (R/W)
\item [SOURCE\_ADDRESS6] Configures whether to which fields of the received frame EMACADDR 6[47:0] compares.\\
0: The DA fields\\
1: The SA fields\\ (R/W)
\item [MASK\_BYTE\_CONTROL6] Configures whether to mask MAC address byte bytes when comparing the received DA or SA with the contents of EMACADDR6.\\
0: Unmask\\
1: Mask\\
The correspondence between the mask control bits and the bytes:\\
Bit[29]: EMACADDR6HIGH\_REG[15:8]\\
Bit[28]: EMACADDR6HIGH\_REG[7:0]\\
Bit[27]: EMACADDR6LOW\_REG[31:24]\\
Bit[24]: EMACADDR6LOW\_REG[7:0]\\
You can filter a group of addresses (known as group address filtering) by masking one or more bytes of the address. (R/W)
\item [MAC\_ADDRESS6\_HI] Configures the upper 16 bits of the seventh 6-byte MAC address [47:32]. (R/W)
\end{reglist}\end{regdesc}
\end{register}

\begin{register}{H}{EMACADDR6LOW\_REG}{0x0074}\label{regdesc:EMACADDR6LOWREG}
\regfield{}{32}{0}{{0x0FFFFFFFF}}%
\reglabel{Reset}\regnewline%
\begin{regdesc}\begin{reglist}
\item [EMACADDR6LOW\_REG] Configures the lower 32 bits of the seventh 6-byte MAC address.\\The content of this field is undefined until loaded by the Application after the initialization process.(R/W)
\end{reglist}\end{regdesc}
\end{register}


\begin{register}{H}{EMACADDR7HIGH\_REG}{0x0078}\label{regdesc:EMACADDR7HIGHREG}
\regfield{ADDRESS\_ENABLE7}{1}{31}{0}%
\regfield{SOURCE\_ADDRESS7}{1}{30}{0}%
\regfield{MASK\_BYTE\_CONTROL7}{6}{24}{{0x00}}%
\regfield{(reserved)}{8}{16}{00000000}%
\regfield{MAC\_ADDRESS7\_HI}{16}{0}{{0x0FFFF}}%
\reglabel{Reset}\regnewline%
\begin{regdesc}\begin{reglist}
\item [ADDRESS\_ENABLE7] Configures whether to enable perfect filtering using the eighth MAC address.\\
0: Disable\\
1: Enable\\ (R/W)
\item [SOURCE\_ADDRESS7] Configures whether to which fields of the received frame EMACADDR7 [47:0] compares.\\
0: The DA fields\\
1: The SA fields\\ (R/W)
\item [MASK\_BYTE\_CONTROL7] Configures whether to mask MAC address byte bytes when comparing the received DA or SA with the contents of EMACADDR7.\\
0: Unmask\\
1: Mask\\
The correspondence between the mask control bits and the bytes:\\
Bit[29]: EMACADDR7HIGH\_REG[15:8]\\
Bit[28]: EMACADDR7HIGH\_REG[7:0]\\
Bit[27]: EMACADDR7LOW\_REG[31:24]\\
Bit[24]: EMACADDR7LOW\_REG[7:0]\\
You can filter a group of addresses (known as group address filtering) by masking one or more bytes of the address. (R/W)
\item [MAC\_ADDRESS7\_HI] Configures the upper 16 bits of the eighth 6-byte MAC address [47:32]. (R/W)
\end{reglist}\end{regdesc}
\end{register}

\begin{register}{H}{EMACADDR7LOW\_REG}{0x007C}\label{regdesc:EMACADDR7LOWREG}
\regfield{}{32}{0}{{0x0FFFFFFFF}}%
\reglabel{Reset}\regnewline%
\begin{regdesc}\begin{reglist}
\item [EMACADDR7LOW\_REG] Configures the lower 32 bits of the eighth 6-byte MAC address.\\The content of this field is undefined until loaded by the Application after the initialization process.(R/W)
\end{reglist}\end{regdesc}
\end{register}

\begin{register}{H}{EMACADDR8HIGH\_REG}{0x0080}\label{regdesc:EMACADDR8HIGHREG}
\regfield{ADDRESS\_ENABLE8}{1}{31}{0}%
\regfield{SOURCE\_ADDRESS8}{1}{30}{0}%
\regfield{MASK\_BYTE\_CONTROL8}{6}{24}{{0x00}}%
\regfield{(reserved)}{8}{16}{00000000}%
\regfield{MAC\_ADDRESS8\_HI}{16}{0}{{0x0FFFF}}%
\reglabel{Reset}\regnewline%
\begin{regdesc}\begin{reglist}
\item [ADDRESS\_ENABLE8] Configures whether to enable perfect filtering using the ninth MAC address.\\
0: Disable\\
1: Enable\\ (R/W)
\item [SOURCE\_ADDRESS8] Configures whether to which fields of the received frame EMACADDR8 [47:0] compares.\\
0: The DA fields\\
1: The SA fields\\ (R/W)
\item [MASK\_BYTE\_CONTROL8] Configures whether to mask MAC address byte bytes when comparing the received DA or SA with the contents of EMACADDR8.\\
0: Unmask\\
1: Mask\\
The correspondence between the mask control bits and the bytes:\\
Bit[29]: EMACADDRH8IGH\_REG[15:8]\\
Bit[28]: EMACADDR8HIGH\_REG[7:0]\\
Bit[27]: EMACADDR8LOW\_REG[31:24]\\
Bit[24]: EMACADDR8LOW\_REG[7:0]\\
You can filter a group of addresses (known as group address filtering) by masking one or more bytes of the address. (R/W)
\item [MAC\_ADDRESS8\_HI] Configures the upper 16 bits of the ninth 6-byte MAC address [47:32]. (R/W)
\end{reglist}\end{regdesc}
\end{register}

\begin{register}{H}{EMACADDR8LOW\_REG}{0x0084}\label{regdesc:EMACADDR8LOWREG}
\regfield{}{32}{0}{{0x0FFFFFFFF}}%
\reglabel{Reset}\regnewline%
\begin{regdesc}\begin{reglist}
\item [EMACADDR8LOW\_REG] Configures the lower 32 bits of the ninth 6-byte MAC address.\\The content of this field is undefined until loaded by the Application after the initialization process. (R/W)
\end{reglist}\end{regdesc}
\end{register}




\begin{register}{H}{EMACCSTATUS\_REG}{0x00D8}\label{regdesc:EMACCSTATUSREG}
\regfield{(reserved)}{15}{17}{000000000000000}%
\regfield{(reserved)}{1}{16}{0}%
\regfield{(reserved)}{11}{5}{00000000000}%
\regfield{JABBER\_TIMEOUT}{1}{4}{0}%
\regfield{(reserved)}{3}{1}{000}%
\regfield{LINK\_MODE}{1}{0}{0}%
\reglabel{Reset}\regnewline%
\begin{regdesc}\begin{reglist}
\item [JABBER\_TIMEOUT]  Represents whether a receive jabber error occurs.\\
0: Not occur\\
1: Occur\\ (RO)
\item [LINK\_MODE]  Represents the current mode of operation of the link.\\
0: Half-duplex\\
1: Full-duplex\\ (RO)
\end{reglist}\end{regdesc}
\end{register}


\begin{register}{H}{EMACWDOGTO\_REG}{0x00DC}\label{regdesc:EMACWDOGTOREG}
\regfield{(reserved)}{15}{17}{000000000000000}%
\regfield{PWDOGEN}{1}{16}{0}%
\regfield{(reserved)}{2}{14}{00}%
\regfield{WDOGTO}{14}{0}{{0x0000}}%
\reglabel{Reset}\regnewline%
\begin{regdesc}\begin{reglist}
\item [PWDOGEN] Configures whether the watchdog timeout is programmble.\\
0: Not programmable. The watchdog timeout for a received frame is controlled by the setting of EMACCONFIG\_REG[23] (WD) and EMACCONFIG\_REG[20] (JE).\\
1: Programmable via the WDOGTO field if EMACCONFIG\_REG[23] (WD) is 0.\\ (R/W)
\item [WDOGTO] Configures the watchdog timeout for a received frame when PWDOGEN is 1 and EMACCONFIG\_REG[23] (WD) is 0.\\
If the length of a received frame exceeds the value of this field, such frame is terminated and declared as an error frame. (R/W)
\end{reglist}\end{regdesc}
\end{register}


\begin{register}{H}{EMACTXVLTCTRL\_REG}{0x0584}\label{regdesc:EMACTXVLANTAGCONTROLREG}
\regfield{(reserved)}{12}{20}{000000000000}%
\regfield{CSVL}{1}{19}{0}%
\regfield{VLP}{1}{18}{0}%
\regfield{VLC}{2}{16}{00}%
\regfield{VLT}{16}{0}{{0x0000}}%
\reglabel{Reset}\regnewline%
\begin{regdesc}\begin{reglist}
\item [CSVL] Configures which type is inserted or replaced in the 13th and 14th bytes of transmitted frames.\\
0: C-VLAN type (0x8100)\\
1: S-VLAN type (0x88A8)\\ (R/W)
\item [VLP] Configures what controls VLAN deletion, insertion, or replacement.\\
0: The control input signal\\
1: VLC\\ (R/W)
\item [VLC] Configures the VLAN tag in transmit frames.\\
0: No VLAN tag deletion, insertion, or replacement.\\
1: VLAN tag deletion. The MAC removes the VLAN type (bytes 13 and 14) and VLAN tag (bytes 15 and 16) of all transmitted frames with VLAN tags.\\
2: VLAN tag insertion. The MAC inserts VLT in bytes 15 and 16 of the frame after inserting the Type value (0x8100/0x88a8) in bytes 13 and 14. This operation is performed on all transmitted frames, irrespective of whether they already have a VLAN tag.\\
3: VLAN tag replacement. The MAC replaces VLT in bytes 15 and 16 of all VLAN-type transmitted frames (Bytes 13 and 14 are 0x8100/0x88a8). \\ (R/W)
\item [VLT] Configures the value of the VLAN tag to be inserted or replaced.\\
The value must only be changed when the transmit lines are inactive or during the initialization phase. (R/W)
\end{reglist}\end{regdesc}
\end{register}

\begin{register}{H}{EMACTSTPCTRL\_REG}{0x0700}\label{regdesc:EMACTIMESTAMPCONTROLREG}
\regfield{(reserved)}{13}{19}{0000000000000}%
\regfield{TSENMACADDR}{1}{18}{0}%
\regfield{SNAPTYPSEL}{2}{16}{00}%
\regfield{TSMSTRENA}{1}{15}{0}%
\regfield{TSEVNTENA}{1}{14}{0}%
\regfield{TSIPV4ENA}{1}{13}{1}%
\regfield{TSIPV6ENA}{1}{12}{0}%
\regfield{TSIPENA}{1}{11}{0}%
\regfield{TSVER2ENA}{1}{10}{0}%
\regfield{TSCTRLSSR}{1}{9}{0}%
\regfield{TSENALL}{1}{8}{0}%
\regfield{(reserved)}{2}{6}{00}%
\regfield{TSADDREG}{1}{5}{0}%
\regfield{TSTRIG}{1}{4}{0}%
\regfield{TSUPDT}{1}{3}{0}%
\regfield{TSINIT}{1}{2}{0}%
\regfield{TSCFUPDT}{1}{1}{0}%
\regfield{TSENA}{1}{0}{0}%
\reglabel{Reset}\regnewline%
\begin{regdesc}\begin{reglist}
\item [TSENMACADDR] Configures whether to enable the DA MAC address (that matches any MAC Address register) for PTP frame filtering when PTP is directly sent over Ethernet.\\
0: Disable\\
1: Enable\\ (R/W)
\item [SNAPTYPSEL] Configures the set of PTP packet types for which snapshot needs to be taken, in combination with TSMSTRENA and TSEVNTENA.\\
SNAPTYPSEL, TSMSTRENA, TSEVNTENA:\\
00x0: SYNC, Follow\_Up, Delay\_Req, Delay\_Resp\\
0001: SYNC\\
0011: Delay\_Req\\
01x0: SYNC, Follow\_Up, Delay\_Req, Delay\_Resp, Pdelay\_Req, Pdelay\_Resp, Pdelay\_Resp\_Follow\_Up\\
0101: SYNC, Pdelay\_Req, Pdelay\_Resp\\
0111: Delay\_Req, Pdelay\_Req, Pdelay\_Resp\\
10xx: SYNC, Delay\_Req\\
11xx: Pdelay\_Req, Pdelay\_Resp\\
(R/W)
\item [TSMSTRENA] Configures whether to enable snapshot only for messages relevant to the master node.\\
0: Disable. The snapshot is taken for the messages relevant to the slave node\\
1: Enable\\ (R/W)
\item [TSEVNTENA] Configures whether to enable timestamp snapshot only for event messages (SYNC, Delay\_Req, Pdelay\_Req, Pdelay\_Resp).\\
0: Disable. The snapshot is taken for all messages except Announce, Management, and Signaling\\
1: Enable\\ (R/W)
\item [TSIPV4ENA] Configures whether to enable the processing of PTP frames sent over IPv4-UDP.\\
0: Disable. The MAC ignores the PTP transported over UDP-IPv4 packets\\
1: Enable. The MAC receiver processes the PTP packets encapsulated in UDP over IPv4 packets\\ (R/W)

\item [Continued on the next page...]
\end{reglist}\end{regdesc}
\end{register}

\addtocounter{Regfloat}{-1}
\begin{register}{H}{EMACTSTPCTRL\_REG}{0x0700}
\begin{regdesc}\begin{reglist}
\item [Continued from the previous page...]

\item [TSIPV6ENA] Configures whether to enable the processing of PTP frames sent over IPv6-UDP.\\
0: Disable. The MAC ignores the PTP transported over UDP-IPv6 packets\\
1: Enable. The MAC receiver processes the PTP packets encapsulated in UDP over IPv6 packets\\ (R/W)
\item [TSIPENA] Configures whether to enable the processing of PTP sent over Ethernet frames.\\
0: Disable. the MAC ignores the PTP over Ethernet packets\\
1: Enable. The MAC receiver processes the PTP packets encapsulated directly in the Ethernet frames\\ (R/W)
\item [TSVER2ENA] Configures whether to enable PTP packet processing for version 2 format.\\
0: Disable. The PTP packets are processed using the version 1 format\\
1: Enable. The PTP packets are processed using the 1588 version 2 format\\ (R/W)
\item [TSCTRLSSR] Configures the timeout digital or binary rollover.\\
0: The Timestamp Low register rolls over after the value of the sub-second register 0x7FFF\_FFFF. The sub-second increment has to be programmed correctly depending on the PTP reference clock frequency and the value of this bit.\\
1: The Timestamp Low register rolls over after 0x3B9A\_C9FF value (that is, 1 nanosecond accuracy) and increments the timestamp (High) seconds.\\
 (R/W)
\item [TSENALL] Configures whether to enable timestamp for all frames received by the MAC.\\
0: Disable\\
1: Enable\\ (R/W)
\item [TSADDREG] Configures whether to update the content of the EMACTSTPADDEND\_REG register in the PTP block for fine correction.\\
0: Not update\\
1: Update\\
This bit should be zero before setting it. (R/WS/SC)

\item [Continued on the next page...]
\end{reglist}\end{regdesc}
\end{register}

\addtocounter{Regfloat}{-1}
\begin{register}{H}{EMACTSTPCTRL\_REG}{0x0700}
\begin{regdesc}\begin{reglist}
\item [Continued from the previous page...]

\item [TSTRIG] Configures whether to enable the timestamp interrupt trigger.\\
0: Disable\\
1: Enable. The timestamp interrupt is generated when the System Time becomes greater than the value written in the Target Time register\\ (R/WS/SC)
\item [TSUPDT] Configures whether to update (add or subtract) the system time with the value specified in  EMAC2NDUPT\_REG and EMACNSUPT\_REG.\\
0: Not update\\
1: Update\\
This bit should be read as zero before updating it.\\
EMACSYSTIMHIGHWORD2ND\_REG is not updated. (R/WS/SC)
\item [TSINIT] Configures whether to initialize (overwrite) the system time with the value specified in  EMAC2NDUPT\_REG and EMACNSUPT\_REG.\\
0: Not initialize\\
1: Initialize\\
This bit should be read as zero before updating it.\\
EMACSYSTIMHIGHWORD2ND\_REG will be initialized. (R/WS/SC)
\item [TSCFUPDT] Configures the method used to update the system time.\\
0: Coarse method\\
1: Find update method\\ (R/W)
\item [TSENA] Configures whether to add the timestamp for the transmit and receive frames.\\
0: Not add the timestamp, and receive frames and the Timestamp Generator is also suspended\\
1: Add the timestamp\\
You need to initialize the Timestamp (system time) after setting this bit to 1.\\
On the receive side, the MAC processes the 1588 frames only if this bit is set. (R/W)
\end{reglist}\end{regdesc}
\end{register}

\begin{register}{H}{EMACSUB2NDINCR\_REG}{0x0704}\label{regdesc:EMACSUBSECONDINCRREG}
\regfield{(reserved)}{24}{8}{000000000000000000000000}%
\regfield{SSINC}{8}{0}{{0x00}}%
\reglabel{Reset}\regnewline%
\begin{regdesc}\begin{reglist}
\item [SSINC] Configures the value accumulated every clock cycle (PTP clock) with the contents of the sub-second register.\\
For example, when PTP clock is 50 MHz (the period is 20 ns), you should program 20 (0x14) when the System Time-Nanoseconds register has an accuracy of 1 ns (EMACTSTPCTRL\_REG[9], i.e., the TSCTRLSSR bit is set). When TSCTRLSSR is cleared, the Nanoseconds register has a resolution of ~0.465 ns. In this case, you should program a value of 43 (0x2B) that is derived by 20 ns/0.465. (R/W)
\end{reglist}\end{regdesc}
\end{register}

\begin{register}{H}{EMACSYSTIM2ND\_REG}{0x0708}\label{regdesc:EMACSYSTEMTIMESREG}
\regfield{}{32}{0}{{0x00000000}}%
\reglabel{Reset}\regnewline%
\begin{regdesc}\begin{reglist}
\item [TSS] Represents the current value of the System Time maintained by the MAC.\\
Measurement unit: second. (RO)
\end{reglist}\end{regdesc}
\end{register}

\begin{register}{H}{EMACSYSTIMNS\_REG}{0x070C}\label{regdesc:EMACSYSTEMTIMENANOSSREG}
\regfield{(reserved)}{1}{31}{0}%
\regfield{TSSS}{31}{0}{{0x00000000}}%
\reglabel{Reset}\regnewline%
\begin{regdesc}\begin{reglist}
\item [TSSS] Represents the sub-second representation of time, with an accuracy of 0.46 ns.\\
When EMACTSTPCTRL\_REG[9] (TSCTRLSSR) is set, each bit represents 1 ns and the maximum value is 0x3B9A\_C9FF. (RO)
\end{reglist}\end{regdesc}
\end{register}

\begin{register}{H}{EMAC2NDUPT\_REG}{0x0710}\label{regdesc:EMACSECONDSUPDATEREG}
\regfield{}{32}{0}{{0x00000000}}%
\reglabel{Reset}\regnewline%
\begin{regdesc}\begin{reglist}
\item [TSSSUPD] Configures the time to be initialized or added to the system time.\\
Measurement unit: second. (R/W)
\end{reglist}\end{regdesc}
\end{register}

\begin{register}{H}{EMACNSUPT\_REG}{0x0714}\label{regdesc:EMACNANOSUPDATEREG}
\regfield{ADDSUB}{1}{31}{0}%
\regfield{TSSSUPD}{31}{0}{{0x00000000}}%
\reglabel{Reset}\regnewline%
\begin{regdesc}\begin{reglist}
\item [ADDSUB]with  Configures whether to subtract or add the system time value the contents of the update register.\\
0: Add\\
1: Subtract\\ (R/W)
\item [TSSSUPD] Configures the sub-second representation of time, with an accuracy of 0.46 ns.\\
When EMACTSTPCTRL\_REG[9] (TSCTRLSSR) is set, each bit represents 1 ns and the programmed value should not exceed 0x3B9A\_C9FF. (R/W)
\end{reglist}\end{regdesc}
\end{register}

\begin{register}{H}{EMACTSTPADDEND\_REG}{0x0718}\label{regdesc:EMACTIMESTAMPADDENDREG}
\regfield{}{32}{0}{{0x00000000}}%
\reglabel{Reset}\regnewline%
\begin{regdesc}\begin{reglist}
\item [TSAR] Configures the 32-bit time value to be added to the Accumulator register to achieve time synchronization. (R/W)
\end{reglist}\end{regdesc}
\end{register}

\begin{register}{H}{EMACTGTIME2ND\_REG}{0x071C}\label{regdesc:EMACTARGETTIMESECONDREG}
\regfield{}{32}{0}{{0x00000000}}%
\reglabel{Reset}\regnewline%
\begin{regdesc}\begin{reglist}
\item [TSTR] Configures the target time.\\
Measurement unit: second.\\
When the system time matches or exceeds the target time, an interrupt is generated (if enabled). (R/W)
\end{reglist}\end{regdesc}
\end{register}


\begin{register}{H}{EMACTGTIMENS\_REG}{0x0720}\label{regdesc:EMACTARGETTIMENANOSREG}
\regfield{(reserved)}{1}{31}{0}%
\regfield{TTSLO}{31}{0}{{0x00000000}}%
\reglabel{Reset}\regnewline%
\begin{regdesc}\begin{reglist}
\item [TTSLO] Configures the signed target time.\\
Measurement unit: nanosecond.\\
When the system time matches or exceeds the target time, an interrupt is generated (if enabled). (R/W)
\end{reglist}\end{regdesc}
\end{register}

\begin{register}{H}{EMACSYSTIMHIGHWORD2ND\_REG}{0x0724}\label{regdesc:EMACSYSTEMTIMEHIGHERWORDSECONDREG}
\regfield{(reserved)}{16}{16}{0000000000000000}%
\regfield{TSHWR}{16}{0}{{0x0000}}%
\reglabel{Reset}\regnewline%
\begin{regdesc}\begin{reglist}
\item [TSHWR] Configures the most significant 16-bits of the timestamp seconds value.\\
The register is directly written to initialize the value.\\
This register is incremented when there is an overflow from the 32 bits of the System Time Seconds register. (R/W/SU)
\end{reglist}\end{regdesc}
\end{register}

\begin{register}{H}{EMACTSTPSTATUS\_REG}{0x0728}\label{regdesc:EMACTIMESTAMPSTATUSREG}
\regfield{(reserved)}{28}{4}{0000000000000000000000000000}%
\regfield{TSTRGTERR}{1}{3}{0}%
\regfield{(reserved)}{1}{2}{0}%
\regfield{TSTARGT}{1}{1}{0}%
\regfield{TSSOVF}{1}{0}{0}%
\reglabel{Reset}\regnewline%
\begin{regdesc}\begin{reglist}
\item [TSTRGTERR] Represents whether the target time, being programmed in EMACTGTIME2ND\_REG and EMACTGTIMENS\_REG, is already elapsed.\\
0: Not elapsed\\
1: Elapsed\\ (R/SS/RC)
\item [TSTARGT] Represents whether the system time is greater or equal to the value specified in EMACTGTIME2ND\_REG and EMACTGTIMENS\_REG.\\
0: Less than the value\\
1: Greater or equal to the value\\ (R/SS/RC)
\item [TSSOVF] Represents whether the second value of the timestamp has overflowed beyond 0xFFFF\_FFFF.\\
0: No overflow\\
1: Overflow\\ (R/SS/RC)
\end{reglist}\end{regdesc}
\end{register}
